% !TeX root=../main.tex
\chapter{نتایج}
%\thispagestyle{empty} 
\label{chap:results}

\section{مقدمه}
این فصل به ارائه، تحلیل و تفسیر نتایج تجربی حاصل از پیاده‌سازی و ارزیابی
روش پیشنهادی در این پژوهش اختصاص دارد.
در حالی‌که فصل سوم به تشریح چارچوب نظری، معماری مدل و راهبردهای آموزش پرداخته است،
تمرکز فصل حاضر بر پاسخ‌گویی عملی و داده‌محور به سوالات تحقیق
و بررسی میزان تحقق اهداف پژوهش از طریق آزمایش‌های کنترل‌شده و قابل تکرار است.
بدین منظور، نتایج حاصل از اجرای مدل‌ها بر روی مجموعه‌ای متنوع از دیتاست‌های
استاندارد ارائه شده و با استفاده از معیارهای کمی و تحلیل‌های آماری
مورد ارزیابی قرار می‌گیرند.

هدف اصلی این پژوهش، ارتقای مدل \lr{Depth Anything V2} از یک مدل
تخمین عمق غیرمتریک\enfootnote{Non-metric Depth}
به مدلی با خروجی عمق متریک پایدار\enfootnote{Metric Depth}
و کاهش حساسیت آن نسبت به تغییر پارامترهای دوربین،
به‌ویژه ماتریس مشخصهٔ دوربین\enfootnote{Camera Intrinsics} بوده است.
ازاین‌رو، در این فصل به‌طور خاص بررسی می‌شود که آیا ادغام صریح
اطلاعات هندسی دوربین در فرآیند تخمین عمق،
می‌تواند ناپایداری مقیاس و وابستگی به نوع دوربین را
که در مدل‌های تک‌دوربینهٔ مرسوم مشاهده می‌شود، کاهش دهد یا خیر.

برای ارزیابی جامع عملکرد \proposedmodel{}،
مجموعه‌ای از آزمایش‌های کمی بر روی دیتاست‌های متنوع شامل
محیط‌های داخلی، خارجی، واقعی و مصنوعی انجام شده است.
در این آزمایش‌ها، \proposedmodel{} با \basemodel{} مقایسه شده
و عملکرد آن با استفاده از معیارهای متداول ارزیابی عمق،
از جمله معیارهای غیرمتریک و متریک، سنجیده می‌شود.
همچنین، به‌منظور بررسی پایداری نتایج،
از آزمون‌های آماری و بازه‌های اطمینان استفاده شده
تا معناداری تفاوت‌های مشاهده‌شده به‌صورت دقیق تحلیل گردد.
نتایج ارائه‌شده در این فصل، مبنای قضاوت نهایی
در خصوص کارایی و اعتبار روش پیشنهادی را فراهم می‌کنند.


\section{سوالات تحقیق و فرضیات تجربی}
\label{sec:research_questions}

در فصل نخست این پژوهش، مجموعه‌ای از سوالات تحقیق با هدف بررسی
امکان دستیابی به عمق متریک پایدار در چارچوب تخمین عمق تک‌دوربینه
و نقش اطلاعات هندسی دوربین در این فرآیند مطرح شد.
در این فصل، این سوالات بر اساس نتایج تجربی به‌دست‌آمده
مورد ارزیابی و پاسخ‌گویی قرار می‌گیرند.

\subsection{بازگویی سوالات تحقیق}
سوالات اصلی تحقیق به‌صورت زیر قابل بازبیان هستند:

\begin{itemize}
    \item آیا \basemodel{}،
    قادر به تولید عمق متریک پایدار در دیتاست‌هایی با دوربین‌ها
    و پارامترهای تصویربرداری متفاوت است؟
    
    \item آیا ادغام صریح ماتریس مشخصهٔ دوربین\enfootnote{Camera Intrinsics}
    در ورودی مدل، می‌تواند وابستگی خروجی عمق به پارامترهای دوربین
    را کاهش داده و پایداری مقیاس فیزیکی را بهبود بخشد؟
    
    \item آیا بهبود عملکرد مدل در معیارهای متریک،
    بدون تخریب ساختار نسبی عمق و معیارهای غیرمتریک
    قابل دستیابی است؟
\end{itemize}

این سوالات، به‌طور مستقیم از محدودیت‌های مدل‌های
تخمین عمق تک‌دوربینهٔ متداول ناشی می‌شوند
که عموماً خروجی‌هایی با ابهام مقیاس\enfootnote{Scale Ambiguity}
تولید می‌کنند و نسبت به تغییر پارامترهای دوربین حساس هستند.

\subsection{فرضیات تجربی}
به‌منظور پاسخ‌گویی نظام‌مند به سوالات فوق،
فرضیات تجربی زیر در این پژوهش تعریف و مورد آزمون قرار گرفته‌اند:

\begin{itemize}
    \item \textbf{فرضیهٔ اول:}
    \basemodel{} در معیارهای متریک،
    نسبت به نوع دیتاست و پارامترهای دوربین حساس بوده
    و تفاوت عملکرد معنادار آماری بین دیتاست‌های مختلف نشان می‌دهد.
    
    \item \textbf{فرضیهٔ دوم:}
    ادغام صریح ماتریس مشخصهٔ دوربین $K$ در \proposedmodel{}،
    منجر به کاهش معنادار خطا در معیارهای متریک
    و بهبود پایداری مقیاس فیزیکی خروجی عمق می‌شود.
    
    \item \textbf{فرضیهٔ سوم:}
    بهبود عملکرد متریک \proposedmodel{}،
    بدون افت محسوس در معیارهای غیرمتریک
    و ساختار نسبی نقشهٔ عمق حاصل می‌شود.
\end{itemize}

\subsection{ارتباط فرضیات با بخش‌های فصل}
برای شفافیت در روند ارزیابی، هر یک از فرضیات فوق
در بخش مشخصی از این فصل مورد بررسی قرار می‌گیرد:

\begin{itemize}
    \item فرضیهٔ اول در بخش
    \ref{sec:quantitative_results}
    و به‌طور خاص در تحلیل عملکرد \basemodel{}
    در دیتاست‌های مختلف بررسی می‌شود.
    
    \item فرضیهٔ دوم در بخش‌های
    مقایسهٔ کمی بین \basemodel{} و \proposedmodel{}
    و همچنین در تحلیل پایداری بین‌دیتاستی
    \proposedmodel{} مورد ارزیابی قرار می‌گیرد.
    
    \item فرضیهٔ سوم در بخش تحلیل معیارهای غیرمتریک
    و نتایج کیفی، از طریق مقایسهٔ بصری نقشه‌های عمق
    بررسی خواهد شد.
\end{itemize}

بدین ترتیب، ساختار فصل به‌گونه‌ای طراحی شده است
که پاسخ‌گویی به سوالات تحقیق و آزمون فرضیات
به‌صورت مرحله‌به‌مرحله و مبتنی بر شواهد تجربی انجام گیرد.

\section{تنظیمات آزمایش و پروتکل ارزیابی}
\label{sec:experimental_setup}

در این بخش، تنظیمات آزمایش‌ها، دیتاست‌های مورد استفاده
و پروتکل ارزیابی مدل‌ها تشریح می‌شود.
هدف از این بخش، ایجاد یک پل منطقی میان چارچوب روش‌شناختی
ارائه‌شده در فصل سوم و نتایج تجربی گزارش‌شده در ادامهٔ این فصل است،
به‌گونه‌ای که تفسیر نتایج بر مبنای شرایط آزمایش‌ها
به‌صورت شفاف و قابل بازتولید انجام گیرد.

\subsection{دیتاست‌های مورد استفاده}
به‌منظور ارزیابی جامع عملکرد مدل‌ها در شرایط متنوع،
از مجموعه‌ای از دیتاست‌های استاندارد و پرکاربرد
در حوزهٔ تخمین عمق تک‌دوربینه استفاده شده است.
این دیتاست‌ها از نظر نوع محیط، منبع داده و پارامترهای دوربین
تنوع بالایی دارند و امکان بررسی پایداری مدل
نسبت به تغییر شرایط تصویربرداری را فراهم می‌کنند.

\begin{itemize}
    \item \textbf{VKITTI:}
    دیتاستی مصنوعی\enfootnote{Synthetic Dataset}
    شامل صحنه‌های بیرونی، تولیدشده بر پایهٔ دیتاست KITTI،
    با عمق مرجع دقیق و چگال.
    
    \item \textbf{\lr{NYUv2:}}
    دیتاستی واقعی از محیط‌های داخلی\enfootnote{Indoor},
    ثبت‌شده با دوربین RGB،
    که عمق مرجع آن با استفاده از حسگرهای فعال به‌دست آمده است.
    
    \item \textbf{DIODE:}
    دیتاستی ترکیبی شامل صحنه‌های داخلی و خارجی،
    با تنوع بالا در نوع دوربین و شرایط تصویربرداری.
    
    \item \textbf{DrivingStereo:}
    دیتاستی واقعی از محیط‌های بیرونی\enfootnote{Outdoor}
    با عمق مرجع متراکم،
    استخراج‌شده از داده‌های استریو با کالیبراسیون دقیق دوربین.
    
    \item \textbf{CityScapes:}
    دیتاستی شهری و واقعی،
    که عمق مرجع آن از طریق تطبیق استریو
    و الگوریتم‌های بازسازی سه‌بعدی به‌دست آمده
    و دارای چگالی و کیفیت ناهمگن در نواحی مختلف تصویر است.
\end{itemize}

این تنوع دیتاست‌ها امکان بررسی رفتار مدل‌ها
در شرایط مختلف تصویربرداری، تفاوت میدان دید،
و تغییر پارامترهای هندسی دوربین را فراهم می‌سازد.

\subsection{نوع محیط‌ها و سناریوهای ارزیابی}
دیتاست‌های مورد استفاده را می‌توان از منظر نوع محیط
به سه دستهٔ اصلی تقسیم کرد:

\begin{itemize}
    \item محیط‌های داخلی (\lr{NYUv2} و بخشی از \lr{DIODE})
    \item محیط‌های خارجی (\lr{DrivingStereo}، \lr{CityScapes} و بخشی از \lr{DIODE})
    \item داده‌های مصنوعی (VKITTI)
\end{itemize}

این تقسیم‌بندی امکان تحلیل پایداری مدل‌ها
در سناریوهای مختلف، از فضاهای بسته با مقیاس محدود
تا صحنه‌های باز شهری با عمق‌های بزرگ،
را فراهم می‌کند.

\subsection{معیارهای ارزیابی}
برای ارزیابی عملکرد مدل‌ها،
از ترکیبی از معیارهای غیرمتریک و متریک استفاده شده است
تا هم ساختار نسبی عمق و هم مقیاس فیزیکی خروجی
مورد بررسی قرار گیرد.

\subsubsection{معیارهای غیرمتریک}
معیارهای غیرمتریک\enfootnote{Non-metric Metrics}
برای ارزیابی صحت نسبی عمق و ترتیب اجسام در صحنه به‌کار می‌روند
و نسبت به مقیاس فیزیکی مطلق حساس نیستند.
در این پژوهش، معیار AbsRel
به‌عنوان نمایندهٔ این دسته استفاده شده است.

\subsubsection{معیارهای متریک}
معیارهای متریک\enfootnote{Metric Metrics}
برای سنجش دقت عمق بر حسب واحدهای فیزیکی واقعی
(مانند متر) به‌کار می‌روند
و نسبت به خطاهای مقیاس حساس هستند.
در این پژوهش، معیار SiLog
به‌عنوان معیار اصلی ارزیابی عمق متریک استفاده شده
و در برخی تحلیل‌ها، معیار RMSE نیز گزارش می‌شود.

\subsection{آزمون‌های آماری}
به‌منظور بررسی معناداری تفاوت‌های مشاهده‌شده
بین مدل‌ها و دیتاست‌ها،
از آزمون‌های آماری استفاده شده است.
در مقایسه‌های بین‌دیتاستی این فصل، آزمون t دو نمونه‌ای\enfootnote{Two-sample t-test}
برای مقایسهٔ میانگین خطاها به‌کار رفته
و سطح معناداری برابر با $0.05$ در نظر گرفته شده است.

در مطالعهٔ کنترل‌شدهٔ بخش~\ref{sec:controlled_nyu_study}
که در آن مدل‌های مختلف بر روی مجموعهٔ آزمون یکسانی ارزیابی می‌شوند،
از آزمون‌های \emph{زوجی} در سطح تک‌تصویر استفاده شده است:
برای هر تصویر آزمون، خطای دو مدل به‌صورت زوجی محاسبه شده
و معناداری تفاوت میانگین‌ها با آزمون t زوجی\enfootnote{Paired t-test}
سنجیده می‌شود.
از آنجا که توزیع خطاهای تک‌تصویری لزوماً نرمال نیست،
آزمون ناپارامتری علامت‌دار ویلکاکسون\enfootnote{Wilcoxon Signed-rank Test}
نیز به‌عنوان سنجهٔ استحکام در کنار آن گزارش می‌شود.
طرح زوجی، واریانس ناشی از دشواری متفاوت تصاویر را حذف کرده
و توان آماری آزمون را به‌طور قابل توجهی افزایش می‌دهد.

علاوه بر این، به‌منظور برآورد پایداری نتایج
و کاهش وابستگی به فرضیات توزیعی،
از بازنمونه‌گیری بوت‌استرپ\enfootnote{Bootstrap Resampling}
(با ده هزار بازنمونه بر روی تفاضل‌های زوجی خطا)
برای محاسبهٔ بازه‌های اطمینان $95\%$ استفاده شده است.
ترکیب این رویکردهای آماری
امکان تحلیل دقیق‌تر تفاوت عملکرد مدل‌ها
و اطمینان از اعتبار نتایج تجربی را فراهم می‌کند.

\section{نتایج کمی: مقایسهٔ عددی مدل‌ها}
\label{sec:quantitative_results}

در این بخش، نتایج کمی حاصل از ارزیابی مدل‌ها ارائه و تحلیل می‌شود.
هدف اصلی، بررسی پایداری \basemodel{}
در دیتاست‌های مختلف و مقایسهٔ آن با \proposedmodel{}
از منظر معیارهای متریک و غیرمتریک است.
تمامی نتایج گزارش‌شده در این بخش
بر اساس اجرای یکسان پروتکل ارزیابی
و با استفاده از آزمون‌های آماری ارائه‌شده در
بخش~\ref{sec:experimental_setup}
به‌دست آمده‌اند.

لازم به توضیح است که آزمایش‌های این بخش، مرحلهٔ نخست مطالعات تجربی این پژوهش را تشکیل می‌دهند و پیش از اصلاح تفکیک مجموعه‌های آموزش و آزمون \lr{NYUv2} (بخش ۳--۷) انجام شده‌اند؛ از این رو، مقادیر مطلق گزارش‌شده برای \lr{NYUv2} در این بخش باید صرفاً به‌عنوان مقایسهٔ نسبی میان مدل‌ها تفسیر شوند. مجموعهٔ نهایی و بدون نشت آزمایش‌ها، شامل مقایسهٔ کنترل‌شده با \basemodel{} و \controlmodel{} هم‌داده و تحلیل آماری زوجی کامل، در بخش~\ref{sec:controlled_nyu_study} ارائه می‌شود.

\subsection{عملکرد \basemodel{} در دیتاست‌های مختلف}
\label{subsec:basic_model_cross_dataset}

در گام نخست، عملکرد \basemodel{}
در دیتاست‌های مختلف مورد بررسی قرار می‌گیرد
تا میزان پایداری آن نسبت به نوع دیتاست
و پارامترهای دوربین ارزیابی شود.
برای این منظور، معیار SiLog
به‌عنوان معیار متریک اصلی استفاده شده است.

نتایج مقایسهٔ بین‌دیتاستی \basemodel{}
در جدول~\ref{tab:basic_model_silog}
ارائه شده است.

\begin{table}[h]
\centering
\caption{مقایسهٔ معیار SiLog \publishedbasemodel{} در دیتاست‌های مختلف}
\label{tab:basic_model_silog}
\begin{tabular}{lccc}
\hline
مقایسهٔ دیتاست‌ها & اختلاف میانگین & $p$-value & نتیجه \\ \hline
\lr{VKITTI vs DrivingStereo} & $0.039$ & $<0.05$ & معنادار \\
\lr{NYUv2 vs DrivingStereo} & $0.101$ & $<0.05$ & معنادار \\
\lr{NYUv2 vs DIODE} & $0.450$ & $<0.05$ & معنادار \\ \hline
\end{tabular}
\end{table}

همان‌طور که مشاهده می‌شود،
\basemodel{} در تمامی مقایسه‌های انجام‌شده
تفاوت عملکرد معنادار آماری بین دیتاست‌ها نشان می‌دهد
($p < 0.05$).
این رفتار بیانگر آن است که خروجی عمق \basemodel{}
به نوع دیتاست و در نتیجه به
پارامترهای دوربین و شرایط تصویربرداری وابسته است.
در نتیجه، می‌توان نتیجه گرفت که
\basemodel{}
قادر به تولید عمق متریک پایدار
در شرایط مختلف نیست.

این ناپایداری، حتی در مواردی که
مدل از نظر ساختار نسبی عمق عملکرد مناسبی دارد،
در معیارهای متریک به‌صورت اختلاف‌های قابل توجه
در مقدار خطا ظاهر می‌شود.

\subsection{مقایسهٔ \basemodel{} و \proposedmodel{}}
\label{subsec:da2_vs_revised}

در ادامه، عملکرد \basemodel{}
با \proposedmodel{} که در آن
ماتریس مشخصهٔ دوربین $K$
به‌صورت صریح در فرآیند تخمین عمق ادغام شده است،
مقایسه می‌شود.
تمرکز اصلی این مقایسه
بر معیارهای متریک و به‌ویژه SiLog است.

جدول~\ref{tab:da2_revised_comparison}
نتایج مقایسهٔ مستقیم دو مدل
در دیتاست‌های مختلف را نشان می‌دهد.

\begin{table}[h]
\centering
\caption{مقایسهٔ عملکرد \publishedbasemodel{} و \proposedmodel{} بر اساس معیار SiLog}
\label{tab:da2_revised_comparison}
\begin{adjustbox}{max width=\textwidth}
\begin{tabular}{lcccc}
\hline
دیتاست & \tablepublishedmodel{} & \tableproposedmodel{} & بهبود (\%) & $p$-value \\ \hline
\lr{NYUv2} & $1.284$ & $\mathbf{0.329}$ & $74.38$ & $\ll 0.05$ \\
DIODE & $0.833$ & $\mathbf{0.377}$ & $54.78$ & $\ll 0.05$ \\
CityScapes & $1.835$ & $\mathbf{0.662}$ & $63.90$ & $\ll 0.05$ \\
DrivingStereo & $1.399$ & $\mathbf{0.414}$ & $70.42$ & $\ll 0.05$ \\ \hline
\end{tabular}
\end{adjustbox}
\end{table}

\begin{figure}[htbp]
    \centering
    \begin{tikzpicture}
    \begin{axis}[
        ybar,
        enlargelimits=0.15,
        legend style={at={(0.5,-0.2)}, anchor=north, legend columns=-1},
        ylabel={خطای SiLog (کمتر بهتر است)},
        symbolic x coords={NYUv2, DIODE, CityScapes, DrivingStereo},
        xtick=data,
        nodes near coords,
        nodes near coords align={vertical},
        width=12cm,
        height=7cm,
        bar width=20pt,
        ymin=0
        ]

    \addplot[fill=blue!40] coordinates {(NYUv2,1.284) (DIODE,0.833) (CityScapes,1.835) (DrivingStereo,1.399)};
    \addplot[fill=green!40] coordinates {(NYUv2,0.329) (DIODE,0.377) (CityScapes,0.662) (DrivingStereo,0.414)};
    \legend{\publishedbasemodel{}, \proposedmodel{}}
    \end{axis}
    \end{tikzpicture}
    \caption{مقایسه عملکرد متریک (معیار SiLog) بین \basemodel{} و \proposedmodel{} در دیتاست‌های مختلف. همان‌طور که مشاهده می‌شود \proposedmodel{} در تمامی موارد خطای بسیار کمتری دارد.}
    \label{fig:cross_dataset_silog}
\end{figure}

نتایج نشان می‌دهد که \proposedmodel{}
در تمامی دیتاست‌ها
کاهش چشمگیر و معنادار آماری
در مقدار خطای SiLog نسبت به \basemodel{} دارد.
درصد بهبود مشاهده‌شده
بین حدود ۵۵ تا ۷۴ درصد متغیر بوده
که بیانگر تأثیر قابل توجه
ادغام اطلاعات هندسی دوربین
بر تخمین عمق متریک است.

این بهبود یکنواخت در دیتاست‌های داخلی،
خارجی و مصنوعی نشان می‌دهد که
\proposedmodel{} قادر است
وابستگی به توزیع داده و نوع دوربین
را به‌طور مؤثری کاهش دهد.

\subsection{تحلیل پایداری بین‌دیتاستی \proposedmodel{}}
\label{subsec:revised_cross_dataset}

در این بخش، پایداری \proposedmodel{}
نسبت به تغییر دیتاست و دوربین بررسی می‌شود.
برای این منظور، عملکرد \proposedmodel{}
بین جفت دیتاست‌های مختلف
با استفاده از آزمون t
و بازه‌های اطمینان بوت‌استرپ
مقایسه شده است.

نتایج نشان می‌دهد که
برخلاف \basemodel{}،
تفاوت عملکرد \proposedmodel{}
در بسیاری از مقایسه‌های بین‌دیتاستی
از نظر آماری معنادار نیست
($p > 0.05$)
و بازه‌های اطمینان محاسبه‌شده
شامل مقدار صفر هستند.

این رفتار بیانگر آن است که
\proposedmodel{} توانسته است
وابستگی خروجی عمق به پارامترهای دوربین
را تا حد زیادی کاهش دهد
و به عمق متریک پایدار دست یابد.
بدین ترتیب، می‌توان نتیجه گرفت که
ادغام صریح ماتریس مشخصهٔ دوربین
نقش کلیدی در بهبود پایداری مقیاس
در تخمین عمق تک‌دوربینه ایفا می‌کند
و پاسخ مثبت و مستقیمی
به سوال اصلی پژوهش
در خصوص امکان دستیابی به عمق متریک پایدار ارائه می‌دهد.

\section{مطالعهٔ کنترل‌شده بر روی پروتکل استاندارد \lr{NYUv2}}
\label{sec:controlled_nyu_study}

در مرحلهٔ دوم مطالعات تجربی این پژوهش، به‌منظور ارائهٔ مقایسه‌ای کاملاً منصفانه، بدون نشت داده و قابل انطباق با ادبیات موضوع، مجموعه‌ای از آزمایش‌های کنترل‌شده بر روی پروتکل استاندارد ارزیابی \lr{NYUv2} طراحی و اجرا شد. در این پروتکل که در اغلب پژوهش‌های مرجع از جمله \lr{Depth Anything V2} به‌کار می‌رود \cite{Eigen2014Depth, DepthAnythingV2}، ارزیابی بر روی زیرمجموعهٔ آزمون ۶۵۴ تصویری \lr{Eigen} با نقشه‌های عمق خام حسگر به‌عنوان مرجع، برش استاندارد \lr{Eigen}، و محدودهٔ عمق \fadecimal{۰٫۰۰۱} تا ۱۰ متر انجام می‌شود. برای مدل‌های نسبی، پیش از محاسبهٔ معیارها، تراز مقیاس و جابجایی به‌ازای هر تصویر در فضای عمق معکوس اعمال می‌گردد و سپس پیش‌بینی تراز‌شده به بازهٔ معتبر ارزیابی محدود می‌شود؛ برای مدل‌های متریک، خروجی مدل بدون هیچ‌گونه ترازی مستقیماً با مرجع مقایسه می‌شود. در بازتولید \lr{KITTI} نیز مطابق جدول ۲ مقالهٔ مرجع از ۶۵۲ تصویر تقسیم \lr{Eigen} و حالت بدون برش استفاده شد.

برای هر اندازهٔ انکودر، چهار مدل در شرایط یکسان مقایسه می‌شوند:
\begin{enumerate}
    \item مقادیر گزارش‌شده در مقالهٔ مرجع \lr{Depth Anything V2} (جدول ۲ مقاله)؛
    \item \publishedbasemodel{}، ارزیابی‌شده با زنجیرهٔ ارزیابی این پژوهش؛
    \item \emph{\controlmodel{}}: معماری پایه بدون پارامترهای دوربین، فاین‌تیون‌شده بر روی همان دادهٔ آموزشی بدون نشت، با ابرپارامترها و برنامهٔ آموزشی کاملاً یکسان با \proposedmodel{}؛
    \item \proposedmodel{}: همان تنظیمات، به‌علاوهٔ ورودی صریح ماتریس مشخصهٔ دوربین.
\end{enumerate}
مقایسهٔ (۴) با (۲) میزان بهبود نسبت به \publishedbasemodel{} را می‌سنجد و مقایسهٔ (۴) با (۳) --- که تنها در دسترسی به $K$ تفاوت دارند --- سهم خالص پارامترهای درونی دوربین را جدا می‌کند.

\subsection{ممیزی نقاط وارسی و مسیر رسیدن به آزمایش نهایی}
\label{subsec:checkpoint_audit}

هدف این بخش گزارش رتبه یا عملکرد نهایی مدل‌ها نیست. پرسش این بخش آن است که
\emph{چرا پروتکل نهایی پژوهش به شکل فعلی طراحی شد؟} برای پاسخ به این پرسش، هر
آزمایش به‌صورت یک زنجیرهٔ سه‌مرحله‌ای گزارش می‌شود: نخست یک مشکل یا الگوی مشکوک
مشاهده شد، سپس معنای آن بررسی شد و در پایان، بر پایهٔ آن یک تصمیم برای آزمایش
بعدی گرفته شد. بنابراین، هر سطر جدول~\ref{tab:experiment_journey} باید از ستون
«مشاهده» به ستون «برداشت» و سپس به ستون «اقدام بعدی» خوانده شود.

نخست، حدود بیست \emph{نقطهٔ وارسی} (وزن‌های ذخیره‌شدهٔ مدل در حین یا پایان
آموزش) ممیزی شدند. این نقاط وارسی اندازه‌های \lr{ViT-S} و \lr{ViT-L}، حالت‌های
عمق نسبی و متریک، و ترکیب‌های مختلف داده را پوشش می‌دادند. در بیشتر اجراها، خطای
AbsRel روی دادهٔ اعتبارسنجی درون‌دامنه‌ای بین \fadecimal{۰٫۰۵} و \fadecimal{۰٫۰۷} بود؛ در این معیار عدد
کمتر بهتر است. بااین‌حال، دو اجرا در تمام دوره‌ها دقیقاً روی \fadecimal{۰٫۳۶۹۸} باقی مانده
بودند و دو نقطهٔ وارسی نسبی \lr{ViT-L} نیز برای همهٔ تصاویر خروجی صفر یا نامعتبر
تولید می‌کردند. ثابت‌ماندن یک معیار در تمام دوره‌ها، همراه با خروجی نامعتبر، یک
نتیجهٔ ضعیف معمولی نیست؛ نشانه‌ای از فروریزش آموزش یا ناسازگاری در ذخیره و
بازیابی مدل است. ازاین‌رو، پایین‌بودن خطای اعتبارسنجی به‌تنهایی برای سالم‌بودن
مدل یا توانایی تعمیم آن کافی در نظر گرفته نشد و هر ادعا با ارزیابی مستقل و
بررسی مستقیم آماره‌های خروجی کنترل شد.

\begingroup
\small
\renewcommand{\arraystretch}{1.25}
\begin{longtable}{@{}>{\raggedleft\arraybackslash}p{0.14\textwidth}>{\raggedleft\arraybackslash}p{0.25\textwidth}>{\raggedleft\arraybackslash}p{0.25\textwidth}>{\raggedleft\arraybackslash}p{0.25\textwidth}@{}}
\caption{مسیر تصمیم‌گیری از نقاط وارسی اولیه تا آزمایش نهایی آگاه از دوربین. این جدول رتبه‌بندی مدل‌ها نیست؛ هر سطر نشان می‌دهد یک مشاهده چگونه به یک برداشت و سپس به تصمیم آزمایشی بعدی منجر شد.}
\label{tab:experiment_journey}\\
\toprule
گام & چه چیزی مشاهده شد؟ & این مشاهده چه معنایی داشت؟ & اقدام بعدی چه بود؟ \\ \midrule
\endfirsthead
\multicolumn{4}{c}{\textit{ادامهٔ جدول~\ref{tab:experiment_journey}}}\\[0.5ex]
\toprule
گام & چه چیزی مشاهده شد؟ & این مشاهده چه معنایی داشت؟ & اقدام بعدی چه بود؟ \\ \midrule
\endhead
\midrule
\multicolumn{4}{l}{\textit{ادامه در صفحهٔ بعد}}\\
\endfoot
\bottomrule
\endlastfoot
۱. ممیزی اولیه
& چند نقطهٔ وارسی خروجی صفر یا نامعتبر داشتند؛ بااین‌حال، منحنی اعتبارسنجی بعضی از آن‌ها عادی به نظر می‌رسید.
& یک معیار اعتبارسنجی به‌تنهایی نمی‌تواند سلامت مدل را تضمین کند.
& خط مبنای منتشرشده بازتولید و خروجی همهٔ مدل‌ها مستقل از فرایند آموزش بررسی شد. \\

۲. ارزیابی صفرنمونه
& پس از فاین‌تیون روی چند مجموعه‌دادهٔ محدود، خطای \lr{KITTI} تقریباً از \fadecimal{۰٫۰۸۳} به \fadecimal{۰٫۲۱۴} افزایش یافت؛ این افت در مدل‌های \tableproposedmodel{} و \tablecontrolmodel{} تقریباً یکسان بود.
& فاین‌تیون محدود، توان تعمیم مدل بنیادین را تضعیف کرده بود و افزودن $K$ علت این افت نبود.
& مسیر ارزیابی صفرنمونه کنار گذاشته شد و مطالعه به پروتکل استاندارد درون‌دامنه‌ای منتقل شد. \\

۳. پاک‌سازی داده
& ۵۲۷ تصویر متعلق به بخش آزمون \lr{NYUv2} در دادهٔ آموزش اولیه پیدا شد.
& نتیجهٔ آن آموزش به نشت داده آلوده بود و نمی‌توانست شاهد معتبری برای تعمیم باشد.
& تمام تصاویر تقسیم آزمون \lr{Eigen} پیش از ساخت مجموعه‌های آموزش و اعتبارسنجی حذف شدند. \\

۴. آزمایش نسبی \lr{ViT-S}
& پس از حذف نشت، مدل \tableproposedmodel{} از \tablepublishedmodel{} و \tablecontrolmodel{} بهتر شد (نسخهٔ نسبی، \lr{ViT-S}).
& نخستین شاهد معتبر به دست آمد؛ اما هنوز معلوم نبود نتیجه در مدل بزرگ‌تر نیز تکرار می‌شود.
& همان آزمایش برای \lr{ViT-L} تکرار شد. \\

۵. فروریزش \lr{ViT-L}
& مدل بزرگ‌تر خروجی تقریباً صفر تولید کرد، درحالی‌که مقدار تابع هزینه ظاهراً غیرعادی نبود.
& نمایش هدفِ آموزش با خروجی مورد انتظار ناسازگار بود؛ بنابراین، منحنی هزینه مشکل را آشکار نمی‌کرد.
& هدف از عمق مستقیم به عمق معکوس تغییر کرد و نوع فضای هدف نیز در نقطهٔ وارسی ذخیره شد. \\

۶. زوم‌ـ‌برش هندسی
& در برش معمولی تصویر، میدان دید تغییر می‌کرد ولی $K$ ثابت می‌ماند و درنتیجه از نظر هندسی نادرست می‌شد. پس از اصلاح $K$، هر دو مدل بهتر شدند، اما مدل \tableproposedmodel{} در ارزیابی نسبی برتری روشنی نداشت.
& تنوع هندسی برای هر دو مدل مفید بود؛ ولی تراز نسبی می‌توانست اثر اطلاعات مقیاس موجود در $K$ را حذف کند.
& ارزیابی به حالت متریک و بدون تراز منتقل شد. \\

۷. ارزیابی متریک
& با حذف تراز نسبی، هرچه تنوع میدان دید در آموزش بیشتر شد، استفادهٔ مدل از اطلاعات دوربین نیز منظم‌تر افزایش یافت.
& اثر $K$ عمدتاً به بازیابی مقیاس متریک مربوط بود و پروتکل نسبی آن را پنهان می‌کرد.
& آزمایشی طراحی شد که مقیاس بدون دانستن کانون دوربین قابل بازیابی نباشد. \\

۸. لرزش کانونی
& در طراحی اولیه، شیب پاسخ مقیاس به تغییر $K$ فقط \fadecimal{۰٫۰۶} بود؛ پس از افزودن لرزش کانونی به دادهٔ آموزش، شیب به \fadecimal{۰٫۹۳}--\fadecimal{۰٫۹۷} رسید. مقدار ایده‌آل این شیب ۱ است.
& شیب \fadecimal{۰٫۰۶} یعنی مدل تقریباً $K$ را نادیده می‌گرفت؛ شیب نزدیک ۱ یعنی مقیاس خروجی تقریباً متناسب با کانون ورودی تغییر می‌کرد.
& لرزش کانونی در طراحی نهایی حفظ شد و ضرورت استفاده از $K$ با آزمایش‌های علّی بررسی شد. \\

۹. فروکاست و حذف $K$
& در آزمایش قبلی، هم لرزش کانونی و هم وزن زیان $\lambda$ تغییر کرده بودند؛ بنابراین، عامل اصلی بهبود هنوز قطعی نبود.
& هم‌بستگی میان $K$ و بهبود عملکرد برای ادعای رابطهٔ علّی کافی نبود.
& با ثابت نگه‌داشتن سایر عوامل نشان داده شد لرزش کانونی عامل بهبود است و جایگزینی یا حذف $K$ صحیح، مقیاس پیش‌بینی را مختل می‌کند. \\
\end{longtable}
\endgroup

جدول را می‌توان در سه فاز خلاصه کرد. در \textbf{فاز نخست، پاک‌سازی شواهد}
(ردیف‌های ۱ تا ۳)، نقاط وارسی ناسالم، تعمیم صفرنمونهٔ تخریب‌شده و نشت داده
شناسایی شدند. خروجی این فاز یک پروتکل درون‌دامنه‌ای بدون نشت بود که در آن نتیجهٔ
\proposedmodel{} و \controlmodel{} قابل مقایسه است.

در \textbf{فاز دوم، اصلاح آموزش و ارزیابی} (ردیف‌های ۴ تا ۷)، ابتدا امکان کسب
بهبود بدون نشت تأیید شد، سپس فروریزش مدل بزرگ‌تر با تغییر هدف به عمق معکوس رفع
گردید. همچنین روشن شد که ارزیابی نسبی، با ترازکردن جداگانهٔ هر تصویر، اطلاعات
مقیاس را حذف می‌کند. به همین دلیل، اثر پارامترهای دوربین باید در حالت متریک و
بدون تراز سنجیده شود.

در \textbf{فاز سوم، آزمون استفادهٔ واقعی از دوربین} (ردیف‌های ۸ و ۹)، لرزش
کانونی باعث شد مقیاس صحنه از ظاهر تصویر به‌تنهایی قابل حدس نباشد. در این شرایط، مدل
برای بازیابی مقیاس ناچار بود از $K$ استفاده کند. افزایش شیب جاروی $K$ از \fadecimal{۰٫۰۶}
به بازهٔ \fadecimal{۰٫۹۳} تا \fadecimal{۰٫۹۷} نشان داد پاسخ مدل از حالت تقریباً بی‌تفاوت به حالتی نزدیک
به پاسخ هندسی ایده‌آل رسیده است. در نهایت، آزمایش‌های فروکاست و حذف $K$ نشان
دادند که این اثر صرفاً هم‌بستگی نیست: وقتی $K$ صحیح حذف یا با مقدار نادرست
جایگزین شد، مقیاس خروجی نیز به‌طور قابل پیش‌بینی مختل شد.

بنابراین، پیام اصلی جدول آن نیست که هر ردیف یک مدل بهتر معرفی می‌کند. پیام جدول
این است که اعتبار نتیجهٔ نهایی به چهار شرط وابسته است: دادهٔ بدون نشت، نقطهٔ
وارسی سالم، ارزیابی متریکِ بدون تراز، و آزمایشی که در آن $K$ برای بازیابی مقیاس
واقعاً ضروری باشد.

\subsection{راستی‌آزمایی زنجیرهٔ ارزیابی: بازتولید نتایج \basemodel{}}
\label{subsec:pipeline_verification}

پیش از هرگونه مقایسه، صحت زنجیرهٔ ارزیابی از طریق بازتولید نتایج گزارش‌شدهٔ مقالهٔ مرجع راستی‌آزمایی شد. جدول~\ref{tab:paper_reproduction} نشان می‌دهد که نتایج به‌دست‌آمده با زنجیرهٔ ارزیابی این پژوهش، با دقت مطلوبی بر مقادیر مقاله منطبق است؛ اختلاف‌های جزئی در محدودهٔ تلورانس متعارف بازتولید این محک قرار دارند.

\begin{table}[h]
\centering
\caption{راستی‌آزمایی زنجیرهٔ ارزیابی: بازتولید نتایج \publishedbasemodel{}، نسخهٔ نسبی، بر روی پروتکل استاندارد. هر معیار در ستونی مستقل گزارش شده است.}
\label{tab:paper_reproduction}
\begin{tabular}{llcccc}
\toprule
\multirow{2}{*}{محک} & \multirow{2}{*}{انکودر}
& \multicolumn{2}{c}{مقالهٔ مرجع}
& \multicolumn{2}{c}{ارزیابی این پژوهش} \\
\cmidrule(lr){3-4}\cmidrule(lr){5-6}
& & AbsRel$\downarrow$ & $\delta_1\uparrow$
& AbsRel$\downarrow$ & $\delta_1\uparrow$ \\ \midrule
\lr{NYUv2} & \lr{ViT-S} & $0.053$ & $0.973$ & $0.051$ & $0.973$ \\
\lr{NYUv2} & \lr{ViT-L} & $0.045$ & $0.979$ & $0.043$ & $0.979$ \\
\lr{KITTI} & \lr{ViT-S} & $0.078$ & $0.936$ & $0.083$ & $0.934$ \\
\lr{KITTI} & \lr{ViT-L} & $0.074$ & $0.946$ & $0.076$ & $0.948$ \\ \bottomrule
\end{tabular}
\end{table}

\subsection{ناکارآمدی فاین‌تیون برای ارزیابی برون‌دامنه‌ای}
\label{subsec:zeroshot_failure}

نخستین یافتهٔ این مرحله، نتیجه‌ای منفی اما تعیین‌کننده در طراحی ادامهٔ آزمایش‌ها بود: فاین‌تیون \basemodel{} بر روی ترکیب مجموعه‌داده‌های محدود (نظیر \lr{VKITTI2}، \lr{Cityscapes}، \lr{DIODE} و \lr{NYUv2})، توانایی تعمیم صفرنمونهٔ مدل را به‌شدت تخریب می‌کند. نتایج دقیق در جدول~\ref{tab:zeroshot_failure} آمده است.

\begin{table}[htbp]
\centering
\caption{ارزیابی بدون نشت نسخه‌های فاین‌تیون‌شدهٔ چندمجموعه‌ای در حالت صفرنمونه، در مقایسه با \publishedbasemodel{}، نسخهٔ نسبی \lr{ViT-S}. هر معیار در ستونی مستقل گزارش شده است.}
\label{tab:zeroshot_failure}
\small
\renewcommand{\arraystretch}{1.12}
\begin{tabularx}{\textwidth}{ll>{\raggedleft\arraybackslash}Xcc}
\toprule
محک & مدل & دادهٔ فاین‌تیون & AbsRel$\downarrow$ & $\delta_1\uparrow$ \\ \midrule
\lr{KITTI} & \tableproposedmodel{} & چندمجموعه‌ای & $0.2139$ & $0.6189$ \\
\lr{KITTI} & \tablepublishedmodel{} & --- & $\mathbf{0.0829}$ & $\mathbf{0.9341}$ \\
\addlinespace
\lr{KITTI} & \tablecontrolmodel{} & چندمجموعه‌ای & $0.2132$ & $0.6179$ \\
\lr{KITTI} & \tablepublishedmodel{} & --- & $\mathbf{0.0829}$ & $\mathbf{0.9341}$ \\
\addlinespace
\lr{NYUv2} & \tableproposedmodel{} & چندمجموعه‌ای، بدون \lr{NYUv2} & $0.0657$ & $0.9578$ \\
\lr{NYUv2} & \tablepublishedmodel{} & --- & $\mathbf{0.0508}$ & $\mathbf{0.9733}$ \\ \bottomrule
\end{tabularx}
\end{table}

افزایش AbsRel \lr{KITTI} از حدود \fadecimal{۰٫۰۸۳} به \fadecimal{۰٫۲۱۴} هم در \proposedmodel{} و هم در \controlmodel{} رخ می‌دهد؛ اختلاف \fadecimal{۰٫۲۱۳۹} و \fadecimal{۰٫۲۱۳۲} در حد نوفه است. بنابراین، پارامترهای دوربین نه علت این افت‌اند و نه راه‌حل آن. این مشاهده نشان می‌دهد که مدل بنیادین آموزش‌دیده بر ده‌ها میلیون تصویر، با فاین‌تیون بر چند هزار تصویر کوچک یا مصنوعی از بازنمایی عمومی خود فاصله می‌گیرد. بر همین اساس، مطابق رویهٔ مقالهٔ مرجع، ادامهٔ مطالعه به فاین‌تیون بر بخش آموزشی یک محک و ارزیابی بر بخش آزمون همان محک منتقل شد.

\subsection{نتایج پروتکل نسبی: مقایسهٔ چهارگانه}
\label{subsec:basic_protocol_results}

جدول‌های \ref{tab:nyu_basic_vits} و \ref{tab:nyu_basic_vitl} نتایج مقایسهٔ چهارگانه را برای دو اندازهٔ انکودر نشان می‌دهند. هر دو مدل فاین‌تیون‌شده بر روی تفکیک بدون نشت (۶۳۶ تصویر آموزش) و با دستورالعمل یکسان (۴۰ دوره، نرخ یادگیری $5\times10^{-6}$، انکودر فریزشده، هدف عمق معکوس مطابق بخش ۳--۷) آموزش دیده‌اند.

\begin{table}[h]
\centering
\caption{مقایسهٔ چهارگانه بر روی مجموعهٔ آزمون استاندارد \lr{NYUv2} (۶۵۴ تصویر، مرجع عمق خام)، انکودر \lr{ViT-S}، پروتکل نسبی با تراز مقیاس و جابجایی.}
\label{tab:nyu_basic_vits}
\begin{adjustbox}{max width=\textwidth}
\begin{tabular}{llccccc}
\hline
مدل & دادهٔ فاین‌تیون & AbsRel$\downarrow$ & $\delta_1\uparrow$ & RMSE$\downarrow$ & SqRel$\downarrow$ & $\log_{10}\downarrow$ \\ \hline
مقالهٔ مرجع (جدول ۲) & --- & $0.053$ & $0.973$ & --- & --- & --- \\
\tablepublishedmodel{} & --- & $0.0508$ & $0.9733$ & $0.2414$ & $0.0245$ & $0.0222$ \\
\tablecontrolmodel{} & \lr{NYUv2} & $0.0507$ & $0.9741$ & $0.2088$ & $0.0186$ & $0.0218$ \\
\textbf{\tableproposedmodel{}} & \lr{NYUv2} & $\mathbf{0.0501}$ & $\mathbf{0.9745}$ & $\mathbf{0.2064}$ & $\mathbf{0.0182}$ & $\mathbf{0.0216}$ \\ \hline
\end{tabular}
\end{adjustbox}
\end{table}

پیش از دستیابی به نتیجهٔ \lr{ViT-L} جدول بعد، دو اجرای \proposedmodel{} و \controlmodel{} با همان دستورالعمل اولیه فروریختند: AbsRel اعتبارسنجی از دورهٔ نخست روی حدود \fadecimal{۰٫۳۴۶} ثابت ماند و خروجی شبکه در تمام پیکسل‌ها دقیقاً صفر شد. کاهش ضریب نرخ یادگیری سر شبکه از ۱۰ به ۱ و افزودن گرم‌سازی نیز مشکل را رفع نکرد. بازتولید ابزارگذاری‌شده نشان داد سهم پیکسل‌های مرده در تنها ۱۸ گام از حدود \fadecimal{۰٫۰۰۱} به ۱ رسید. پس از انتقال هدف به عمق معکوس، یک اجرای آزمایشی ۱۵۰ گامی هیچ پیکسل مرده‌ای نداشت و هزینهٔ آن ۲۰ تا ۵۰ برابر کمتر بود؛ سپس آموزش کامل تکرار شد. بنابراین نتیجهٔ \lr{ViT-L} زیر مربوط به نسخهٔ اصلاح‌شدهٔ فضای هدف است، نه اجرای فروریخته.

\begin{table}[h]
\centering
\caption{مقایسهٔ چهارگانه بر روی مجموعهٔ آزمون استاندارد \lr{NYUv2}، انکودر \lr{ViT-L}، پروتکل نسبی.}
\label{tab:nyu_basic_vitl}
\begin{adjustbox}{max width=\textwidth}
\begin{tabular}{llccc}
\hline
مدل & دادهٔ فاین‌تیون & AbsRel$\downarrow$ & $\delta_1\uparrow$ & RMSE$\downarrow$ \\ \hline
مقالهٔ مرجع (جدول ۲) & --- & $0.045$ & $0.979$ & --- \\
\tablepublishedmodel{} & --- & $0.0425$ & $0.9789$ & $0.2119$ \\
\tablecontrolmodel{} & \lr{NYUv2} & $0.0378$ & $\mathbf{0.9862}$ & $0.1922$ \\
\textbf{\tableproposedmodel{}} & \lr{NYUv2} & $\mathbf{0.0377}$ & $0.9861$ & $\mathbf{0.1913}$ \\ \hline
\end{tabular}
\end{adjustbox}
\end{table}

برای تکمیل گزارش اندازهٔ اثر، جدول~\ref{tab:nyu_basic_effect_sizes} تفاضل زوجی و بازهٔ اطمینان بوت‌استرپ چهار مقایسهٔ اصلی \lr{ViT-S} را ارائه می‌کند. علامت مثبت به‌معنای خطای بیشتر مدل مرجع یا \controlmodel{} نسبت به \proposedmodel{} است.

\begin{table}[htbp]
\centering
\caption{اندازهٔ اثر زوجی \proposedmodel{} \lr{ViT-S} بر ۶۵۴ تصویر \lr{NYUv2}.}
\label{tab:nyu_basic_effect_sizes}
\small
\begin{adjustbox}{max width=\textwidth}
\begin{tabular}{llcc}
\hline
مقایسه & معیار & میانگین تفاضل & بازهٔ اطمینان ۹۵٪ \\ \hline
\tablepublishedmodel{} & RMSE & $+0.0350$ & $[+0.0270,+0.0438]$ \\
\tablepublishedmodel{} & AbsRel & $+0.0007$ & $[-0.0008,+0.0022]$ \\
\tablecontrolmodel{} & RMSE & $+0.0024$ & $[+0.0013,+0.0034]$ \\
\tablecontrolmodel{} & AbsRel & $+0.0006$ & $[+0.0002,+0.0010]$ \\ \hline
\end{tabular}
\end{adjustbox}
\end{table}

برای سنجش معناداری آماری تفاوت‌ها، از آزمون‌های زوجی در سطح تک‌تصویر (آزمون t زوجی و آزمون علامت‌دار ویلکاکسون بر روی $n=654$ زوج) استفاده شده است؛ نتایج در جدول~\ref{tab:nyu_basic_significance} آمده است.

\begin{table}[h]
\centering
\caption{آزمون‌های زوجی معناداری بر روی ۶۵۴ تصویر آزمون \lr{NYUv2} (پروتکل نسبی)؛ مقایسهٔ \proposedmodel{} با \publishedbasemodel{} و \controlmodel{}.}
\label{tab:nyu_basic_significance}
\small
\begin{adjustbox}{max width=\textwidth}
\begin{tabular}{llccc}
\hline
انکودر & مدل مرجع & معیار & $p$ (t زوجی) & $p$ (ویلکاکسون) \\ \hline
\lr{ViT-S} & \tablepublishedmodel{} & AbsRel & $0.40$ & $0.09$ \\
\lr{ViT-S} & \tablepublishedmodel{} & RMSE & $2.8\times10^{-15}$ & $5\times10^{-18}$ \\
\lr{ViT-S} & \tablepublishedmodel{} & $\delta_1$ & $0.34$ & $0.05$ \\
\lr{ViT-S} & \tablecontrolmodel{} & AbsRel & $1.3\times10^{-3}$ & $5.0\times10^{-3}$ \\
\lr{ViT-S} & \tablecontrolmodel{} & RMSE & $1.1\times10^{-5}$ & $3.5\times10^{-5}$ \\
\lr{ViT-S} & \tablecontrolmodel{} & $\delta_1$ & $0.29$ & $0.60$ \\ \hline
\lr{ViT-L} & \tablepublishedmodel{} & AbsRel & $2.5\times10^{-8}$ & $8.9\times10^{-6}$ \\
\lr{ViT-L} & \tablepublishedmodel{} & RMSE & $1.0\times10^{-7}$ & $1.1\times10^{-9}$ \\
\lr{ViT-L} & \tablepublishedmodel{} & $\delta_1$ & $4.1\times10^{-6}$ & $1.2\times10^{-12}$ \\
\lr{ViT-L} & \tablecontrolmodel{} & AbsRel & $0.77$ & $0.18$ \\
\lr{ViT-L} & \tablecontrolmodel{} & RMSE & $0.29$ & $0.50$ \\
\lr{ViT-L} & \tablecontrolmodel{} & $\delta_1$ & $0.81$ & $0.04$ \\ \hline
\end{tabular}
\end{adjustbox}
\end{table}

سه نتیجهٔ اصلی از این جدول‌ها قابل استخراج است. نخست، \proposedmodel{} \lr{ViT-L} در هر سه معیار، به‌صورت معنادار آماری از \publishedbasemodel{} بهتر است (کاهش AbsRel از \fadecimal{۰٫۰۴۲۵} به \fadecimal{۰٫۰۳۷۷} و کاهش \fadecimal{۹٫۷} درصدی RMSE، همگی با $p < 10^{-5}$). دوم، در اندازهٔ \lr{ViT-S}، سهم خالص پارامترهای دوربین در برابر \controlmodel{} هم‌داده، کوچک اما از نظر آماری معنادار است (AbsRel با $p=1.3\times10^{-3}$ و RMSE با $p=1.1\times10^{-5}$). سوم --- و از منظر علمی آموزنده --- در اندازهٔ \lr{ViT-L} تفاوت \proposedmodel{} و \controlmodel{} در پروتکل نسبی از نظر آماری معنادار نیست؛ تحلیل این پدیده در ادامهٔ همین بخش و در بخش بحث ارائه می‌شود.

اندازهٔ اثر \lr{ViT-L} در برابر \publishedbasemodel{} نیز روشن است. برای AbsRel، میانگین تفاضل زوجی و بازهٔ اطمینان ۹۵٪ به‌ترتیب
$+0.0048$ و $[+0.0032,+0.0066]$ است. برای RMSE، مقادیر متناظر برابر
$+0.0205$ و $[+0.0133,+0.0282]$ هستند. در مقابل، آزمون‌های \proposedmodel{} و \controlmodel{} برای AbsRel و RMSE به‌ترتیب $p=0.18$ و $p=0.50$ در ویلکاکسون هستند و تفاوت معناداری میان دو بازو در این پروتکل نشان نمی‌دهند.

\subsection{اثر افزونه‌سازی هندسی و یک نتیجهٔ منفی آموزنده}
\label{subsec:aug_basic_results}

فرضیهٔ محتمل برای عدم مشاهدهٔ سهم پارامترهای دوربین در \lr{ViT-L}، ثابت‌بودن بردار $\mathbf{k}$ در مجموعه‌دادهٔ تک‌دوربینهٔ \lr{NYUv2} بود. برای آزمون این فرضیه، زوج مدل \lr{ViT-L} با افزونه‌سازی هندسی آگاه از دوربین (بخش ۳--۷، بازهٔ میدان دید ۴۱ تا ۶۳ درجه) بازآموزی شد. نتایج (جدول~\ref{tab:nyu_basic_aug}) دو یافته به همراه داشت: نخست، خودِ افزونه‌سازی هندسی برای هر دو مدل سودمندی چشمگیری دارد (کاهش AbsRel \controlmodel{} از \fadecimal{۰٫۰۳۷۸} به \fadecimal{۰٫۰۳۵۹} با $p=1.2\times10^{-13}$ و \proposedmodel{} از \fadecimal{۰٫۰۳۷۷} به \fadecimal{۰٫۰۳۶۲} با $p=1.2\times10^{-8}$) و بهترین نتایج این پژوهش در پروتکل نسبی را رقم می‌زند. دوم، حتی با وجود سیگنال هندسی متنوع و صحیح، \proposedmodel{} در پروتکل نسبی بر \controlmodel{} برتری نیافت (و در AbsRel با اختلاف بسیار جزئی \fadecimal{۰٫۰۰۰۳} ولی معنادار، $p=8.1\times10^{-3}$، پشت سر آن قرار گرفت).

\begin{table}[h]
\centering
\caption{اثر افزونه‌سازی هندسی آگاه از دوربین بر زوج مدل \lr{ViT-L} در پروتکل نسبی \lr{NYUv2}.}
\label{tab:nyu_basic_aug}
\begin{adjustbox}{max width=\textwidth}
\begin{tabular}{lccc}
\hline
مدل & AbsRel$\downarrow$ & $\delta_1\uparrow$ & RMSE$\downarrow$ \\ \hline
\tablecontrolmodel{} & $\mathbf{0.0359}$ & $\mathbf{0.9872}$ & $\mathbf{0.1849}$ \\
\tableproposedmodel{} & $0.0362$ & $0.9869$ & $0.1881$ \\ \hline
\end{tabular}
\end{adjustbox}
\end{table}

این نتیجهٔ منفی، تبیینی اصولی دارد که یکی از یافته‌های مفهومی این پژوهش محسوب می‌شود: در پروتکل ارزیابی نسبی، پیش از محاسبهٔ هر معیار، یک تراز مقیاس و جابجایی به‌ازای هر تصویر میان پیش‌بینی و مرجع اعمال می‌شود. اطلاعات اصلی که میدان دید و فاصلهٔ کانونی در اختیار مدل قرار می‌دهند، دقیقاً همان \emph{مقیاس مطلق صحنه} است؛ کمیتی که سازوکار تراز، آن را از ارزیابی حذف می‌کند. به بیان دیگر، پروتکل نسبی به‌صورت ساختاری نسبت به سهم اصلی پارامترهای درونی دوربین «کور» است. این تحلیل پیش‌بینی می‌کند که سهم پارامترهای دوربین باید در ارزیابی \emph{متریک} (بدون تراز) آشکار شود؛ پیش‌بینی‌ای که در زیربخش بعد به‌صورت تجربی آزموده می‌شود.

\subsection{ارزیابی متریک: آشکارشدن سهم پارامترهای درونی دوربین}
\label{subsec:metric_protocol_results}

در این مرحله، زوج مدل \lr{ViT-L} در حالت متریک (تابع هزینهٔ SiLog بر عمق مطلق، مطابق فصل سوم) و با افزونه‌سازی هندسی آموزش دیده و بر روی مجموعهٔ آزمون استاندارد، \emph{بدون هیچ‌گونه تراز مقیاس}، ارزیابی شد. مرجع مقایسه در این حالت، \publishedbasemodel{}، نسخهٔ متریک (فاین‌تیون‌شده بر \lr{Hypersim} برای صحنه‌های داخلی) است. آزمایش در دو دور با دو شدت افزونه‌سازی انجام شد: دور نخست با بازهٔ میدان دید ۴۱ تا ۶۳ درجه و دور دوم با بازهٔ وسیع‌تر ۳۳ تا ۶۳ درجه.

\begin{table}[h]
\centering
\caption{ارزیابی متریک (عمق مطلق، بدون تراز) بر روی مجموعهٔ آزمون استاندارد \lr{NYUv2}، انکودر \lr{ViT-L}. دور ۱: افزونه‌سازی با میدان دید ۴۱--۶۳ درجه؛ دور ۲: ۳۳--۶۳ درجه.}
\label{tab:nyu_metric_results}
\begin{adjustbox}{max width=\textwidth}
\begin{tabular}{lllccc}
\hline
مدل & دادهٔ آموزش & دور & AbsRel$\downarrow$ & $\delta_1\uparrow$ & RMSE$\downarrow$ \\ \hline
\tablepublishedmodel{} & \lr{Hypersim} & --- & $0.2073$ & $0.6919$ & $0.6063$ \\ \hline
\tablecontrolmodel{} & \lr{NYUv2} & ۱ & $0.0796$ & $0.9583$ & $0.2919$ \\
\textbf{\tableproposedmodel{}} & \lr{NYUv2} & ۱ & $\mathbf{0.0791}$ & $\mathbf{0.9593}$ & $\mathbf{0.2882}$ \\ \hline
\tablecontrolmodel{} & \lr{NYUv2} & ۲ & $0.0826$ & $0.9546$ & $0.2998$ \\
\textbf{\tableproposedmodel{}} & \lr{NYUv2} & ۲ & $\mathbf{0.0817}$ & $\mathbf{0.9569}$ & $\mathbf{0.2944}$ \\ \hline
\end{tabular}
\end{adjustbox}
\end{table}

نتایج جدول~\ref{tab:nyu_metric_results} و آزمون‌های زوجی متناظر (جدول~\ref{tab:nyu_metric_significance}) سه یافتهٔ کلیدی دارند:

\begin{table}[h]
\centering
\caption{آزمون‌های زوجی معناداری در ارزیابی متریک ($n=654$)؛ مقایسهٔ \proposedmodel{} با \publishedbasemodel{}، نسخهٔ متریک، و \controlmodel{}.}
\label{tab:nyu_metric_significance}
\small
\begin{adjustbox}{max width=\textwidth}
\begin{tabular}{llccc}
\hline
دور & مدل مرجع & معیار & $p$ (t زوجی) & $p$ (ویلکاکسون) \\ \hline
۱ & \tablepublishedmodel{} & AbsRel & $4.8\times10^{-133}$ & $8.4\times10^{-96}$ \\
۱ & \tablepublishedmodel{} & RMSE & $6.7\times10^{-100}$ & $5.6\times10^{-87}$ \\
۱ & \tablecontrolmodel{} & AbsRel & $0.26$ & $0.53$ \\
۱ & \tablecontrolmodel{} & RMSE & $4.9\times10^{-3}$ & $9.3\times10^{-3}$ \\ \hline
۲ & \tablecontrolmodel{} & AbsRel & $0.089$ & $\mathbf{0.040}$ \\
۲ & \tablecontrolmodel{} & RMSE & $\mathbf{6.5\times10^{-5}}$ & $1.2\times10^{-4}$ \\ \hline
\end{tabular}
\end{adjustbox}
\end{table}

نخست، هر دو مدل فاین‌تیون‌شده، \publishedbasemodel{}، نسخهٔ متریک را با فاصلهٔ بسیار زیاد پشت سر می‌گذارند (کاهش AbsRel از \fadecimal{۰٫۲۰۷} به حدود \fadecimal{۰٫۰۸}، معادل بهبود \fadecimal{۲٫۶} برابری، با $p$ در مرتبهٔ $10^{-133}$)؛ این نتیجه هم‌راستا با گزارش خود مقالهٔ مرجع است که مقیاس متریک مدل \lr{Hypersim} به صحنه‌های واقعی به‌خوبی منتقل نمی‌شود. دوم، برخلاف پروتکل نسبی، در ارزیابی متریک \proposedmodel{} در \emph{هر دو دور} و در \emph{تمامی معیارها} از \controlmodel{} هم‌داده پیشی می‌گیرد و این برتری در معیار RMSE در هر دو دور معنادار آماری است ($p=4.9\times10^{-3}$ در دور ۱ و $p=6.5\times10^{-5}$ در دور ۲). سوم --- و مهم‌تر --- با افزایش شدت تنوع میدان دید از دور ۱ به دور ۲، فاصلهٔ \proposedmodel{} از \controlmodel{} در همهٔ معیارها افزایش می‌یابد و معناداری AbsRel نیز به آستانهٔ تصمیم می‌رسد (ویلکاکسون $p=0.040$). این الگوی \emph{پاسخ به دوز}\enfootnote{Dose--Response Pattern} --- یعنی رشد سهم پارامترهای دوربین متناسب با میزان ابهام هندسی موجود در داده --- که در دو اجرای مستقل آموزش تکرار شده است، شاهد تجربی مهمی بر سودمندی واقعی ورودی پارامترهای درونی دوربین است. زیربخش بعد نشان می‌دهد که با رساندن ابهام هندسی به حد نهایی خود، این سهم از یک برتری آماریِ مرزی به یک اثر قاطع و همه‌جانبه تبدیل می‌شود.

\subsection{دور سوم --- افزونه‌سازی لرزش کانونی: جهش سهم پارامترهای دوربین}
\label{subsec:focal_jitter_results}

نقطهٔ اوج الگوی پاسخ به دوز، به‌کارگیری افزونه‌سازی لرزش کانونی (بخش~\ref{subsec:focal_jitter}) است که مقیاس مطلق صحنه را بدون دسترسی به $K$ اساساً غیرقابل بازیابی می‌کند. زوج مدل \lr{ViT-L} متریک با دستور آموزشی یکسان با دور ۲، به‌علاوهٔ لرزش کانونی ($m=1.6$)، ضریب $\lambda=0.3$ در هزینهٔ SiLog، و قفل بذر تصادفی بازآموزی شد؛ قفل بذر تضمین می‌کند که هر دو بازو دقیقاً همان دنباله‌ی داده، وارونه‌سازی و ضرایب $\alpha$ را دیده باشند (تابع هزینهٔ گام صفر هر دو اجرا تا آخرین رقم اعشار برابر بود: \fadecimal{۱٫۷۲۵}).

\begin{table}[h]
\centering
\caption{دور ۳ (لرزش کانونی): ارزیابی متریک (عمق مطلق، بدون تراز) بر مجموعهٔ آزمون استاندارد \lr{NYUv2}، انکودر \lr{ViT-L}، به‌همراه آزمون زوجی ویلکاکسون ($n=654$) و بازهٔ اطمینان ۹۵٪ بوت‌استرپ تفاضل زوجی.}
\label{tab:nyu_metric_focal}
\small
\begin{adjustbox}{max width=\textwidth}
\begin{tabular}{lcccc}
\hline
معیار & \tableproposedmodel{} & \tablecontrolmodel{} & تفاضل [بازهٔ اطمینان ۹۵٪] & $p$ (ویلکاکسون) \\ \hline
AbsRel$\downarrow$ & $\mathbf{0.0849}$ & $0.1030$ & $-0.0182\ [-0.0210, -0.0153]$ & $4.0\times10^{-31}$ \\
RMSE$\downarrow$ & $\mathbf{0.3071}$ & $0.3431$ & $-0.0361\ [-0.0432, -0.0289]$ & $1.0\times10^{-24}$ \\
$\delta_1\uparrow$ & $\mathbf{0.9520}$ & $0.9226$ & $+0.0294\ [+0.0224, +0.0367]$ & $2.5\times10^{-24}$ \\
SILog$\downarrow$ & $\mathbf{0.0942}$ & $0.1045$ & $-0.0102\ [-0.0118, -0.0087]$ & $1.5\times10^{-34}$ \\
$\log_{10}\downarrow$ & $\mathbf{0.0363}$ & $0.0427$ & $-0.0065\ [-0.0076, -0.0053]$ & $2.8\times10^{-27}$ \\ \hline
\end{tabular}
\end{adjustbox}
\end{table}

نتایج (جدول~\ref{tab:nyu_metric_focal}) جهشی کیفی نسبت به دورهای قبل نشان می‌دهد: \proposedmodel{} در \emph{تمامی} معیارها با معناداری در مرتبهٔ $p \le 10^{-24}$ از \controlmodel{}ِ هم‌داده و هم‌بذر پیشی می‌گیرد؛ در حالی که بهترین معناداری AbsRel در دور ۲ تنها $p=0.040$ بود. توزیع نسبت مقیاس هر تصویر (میانهٔ پیش‌بینی تقسیم بر میانهٔ مرجع) سازوکار این برتری را آشکار می‌کند: \proposedmodel{} حول \fadecimal{۱٫۰۱۴} متمرکز است، در حالی که نوفهٔ کاهش‌ناپذیر $\alpha$ \controlmodel{} را اریب و پراکنده کرده است (میانهٔ \fadecimal{۱٫۰۸۱}). به بیان دیگر، \controlmodel{} --- که ناچار است به میانگین توزیع $\alpha$ رگرسیون کند --- دقیقاً همان رفتاری را نشان می‌دهد که تحلیل بدساختی مسئله در بخش~\ref{subsec:focal_jitter} پیش‌بینی می‌کرد.

سنجش تعمیم به میدان دیدهای متفاوت نیز همین برتری را تأیید می‌کند: در ارزیابی با برش‌های مرکزی قطعی (زوم \fadecimal{۱٫۰}، \fadecimal{۱٫۴} و \fadecimal{۱٫۸} با به‌روزرسانی دقیق $K$، هر ۶۵۴ تصویر)، AbsRel \proposedmodel{} به‌ترتیب \fadecimal{۰٫۰۹۲}، \fadecimal{۰٫۰۹۷} و \fadecimal{۰٫۱۰۷} و \controlmodel{} \fadecimal{۰٫۱۱۰}، \fadecimal{۰٫۱۱۰} و \fadecimal{۰٫۱۱۵} است --- برتری در تمام سطوح زوم.

دو ملاحظهٔ صادقانه ضروری است. نخست، این برتری قاطع نسبت به کنترل، با بهای اندکی در کیفیت مطلق همراه است: AbsRel \proposedmodel{} دور ۳ (\fadecimal{۰٫۰۸۴۹}) اندکی بالاتر از بهترین مدل خام پژوهش (دور ۱، \fadecimal{۰٫۰۷۹۱}) است؛ ادامهٔ همان مصالحه‌ای که از دور ۱ به دور ۲ نیز مشاهده شد: وظیفهٔ آموزشی دشوارتر، کیفیت مطلق را اندکی می‌کاهد اما شواهد سازوکاری را به‌شدت تقویت می‌کند. دوم، دور ۳ نسبت به دور ۲ دو متغیر را هم‌زمان تغییر می‌دهد (افزونه‌سازی لرزش کانونی و $\lambda$)؛ آزمایش تفکیک این دو عامل در زیربخش~\ref{subsec:focal_ablation} گزارش می‌شود.

\subsection{مطالعهٔ فروکاست: تفکیک اثر لرزش کانونی از ضریب $\lambda$}
\label{subsec:focal_ablation}

برای انتساب دقیق اثر مشاهده‌شده در دور ۳، زوج مدل چهارمی با \emph{تنها یک} متغیر تغییریافته نسبت به دور ۲ آموزش داده شد: لرزش کانونی فعال ($m=1.6$) اما $\lambda=0.5$ بدون تغییر (به‌همراه همان قفل بذر تصادفی).

\begin{table}[h]
\centering
\caption{مطالعهٔ فروکاست (لرزش کانونی با $\lambda=0.5$ بدون تغییر): ارزیابی متریک بر مجموعهٔ آزمون استاندارد \lr{NYUv2} ($n=654$).}
\label{tab:nyu_metric_focal_ablation}
\small
\begin{adjustbox}{max width=\textwidth}
\begin{tabular}{lcccc}
\hline
معیار & \tableproposedmodel{} & \tablecontrolmodel{} & تفاضل [بازهٔ اطمینان ۹۵٪] & $p$ (ویلکاکسون) \\ \hline
AbsRel$\downarrow$ & $\mathbf{0.0834}$ & $0.1065$ & $-0.0231\ [-0.0260, -0.0202]$ & $3.4\times10^{-45}$ \\
RMSE$\downarrow$ & $\mathbf{0.2987}$ & $0.3497$ & $-0.0510\ [-0.0583, -0.0437]$ & $7.5\times10^{-40}$ \\
$\delta_1\uparrow$ & $\mathbf{0.9554}$ & $0.9173$ & $+0.0381\ [+0.0310, +0.0453]$ & $1.2\times10^{-39}$ \\
SILog$\downarrow$ & $\mathbf{0.0918}$ & $0.1056$ & $-0.0139\ [-0.0155, -0.0123]$ & $2.2\times10^{-50}$ \\ \hline
\end{tabular}
\end{adjustbox}
\end{table}

نتیجهٔ فروکاست (جدول~\ref{tab:nyu_metric_focal_ablation}) از دو جهت تعیین‌کننده است. نخست، \emph{کل اثر دور ۳ با لرزش کانونی به‌تنهایی بازتولید --- و حتی تقویت --- می‌شود}: برتری \proposedmodel{} در تمامی معیارها به مرتبهٔ $p \le 10^{-39}$ می‌رسد و شیب جاروی $K$ به $0.97 \pm 0.03$ --- نزدیک‌تر به مقدار ایده‌آل ۱ نسبت به دور ۳. بنابراین عامل علّی، خودِ افزونه‌سازی لرزش کانونی است و کاهش $\lambda$ نه‌تنها ضروری نیست، بلکه اندکی هم زیان‌بار بوده است (AbsRel \proposedmodel{} با $\lambda=0.5$: \fadecimal{۰٫۰۸۳۴} در برابر \fadecimal{۰٫۰۸۴۹} با $\lambda=0.3$). دوم، این چهارمین زوج آموزشی مستقلی است که برتری \proposedmodel{} را تکرار می‌کند و در کنار دورهای ۱ تا ۳، الگوی پاسخ به دوز را به‌صورت یکنواخت کامل می‌سازد: هرچه بازیابی مقیاس بدون $K$ دشوارتر، سهم پارامترهای دوربین بزرگ‌تر و معنادارتر.

\subsection{جمع‌بندی مطالعهٔ کنترل‌شده و گزارش کوشش‌های تجربی}
\label{subsec:controlled_study_summary}

برای شفافیت علمی، گسترهٔ کوشش‌های تجربی این مرحله به‌اختصار گزارش می‌شود. در مجموع بیش از بیست اجرای آموزش مستقل (شامل زوج‌های پیشنهادی/کنترل در دو اندازهٔ انکودر، دو نوع مدل نسبی و متریک، و دو شدت افزونه‌سازی) و ممیزی کامل نقاط وارسی پیشین پروژه انجام شد. در جریان این ممیزی، تعدادی از نقاط وارسی قدیمی \lr{ViT-L} به‌دلیل فروریزش خاموش توصیف‌شده در بخش ۳--۷ عملاً بلااستفاده تشخیص داده شدند و سه اشکال زیرساختی (نشت تقسیم‌بندی \lr{NYUv2}، ناهم‌خوانی فضای هدف آموزش، و ناهم‌خوانی معماری سر شبکه در ارزیابی متریک) شناسایی و رفع گردید (بخش ۳--۹). رویکردهای آزموده‌شده و کنارگذاشته‌شده نیز مستند شده‌اند: فاین‌تیون چندمجموعه‌ای برای ارزیابی صفرنمونه (به‌دلیل تخریب تعمیم)، و کاهش نرخ یادگیری سر شبکه به‌عنوان راه‌حل فروریزش (که مؤثر نبود و نشان داد ریشهٔ مشکل در فضای هدف است، نه نرخ یادگیری).

خلاصهٔ یافته‌های این مطالعهٔ کنترل‌شده چنین است: (۱) \proposedmodel{} در هر دو اندازهٔ انکودر و در هر دو پروتکل، از \publishedbasemodel{} به‌صورت معنادار بهتر است؛ (۲) سهم خالص پارامترهای درونی دوربین در پروتکل نسبی تنها در اندازهٔ کوچک معنادار است، زیرا سازوکار تراز، اطلاعات مقیاس را حذف می‌کند؛ (۳) در ارزیابی متریک --- که مقیاس مطلق اهمیت دارد --- سهم پارامترهای دوربین آشکار، معنادار و متناسب با میزان تنوع هندسی داده است؛ و (۴) هنگامی که افزونه‌سازی لرزش کانونی مقیاس را بدون $K$ غیرقابل بازیابی می‌کند، این سهم به برتری قاطع در تمامی معیارها ($p \le 10^{-24}$) می‌رسد.

\section{تحلیل حساسیت نسبت به پارامترهای دوربین}
\label{sec:camera_sensitivity}

یکی از اهداف اصلی این پژوهش،
کاهش حساسیت تخمین عمق نسبت به پارامترهای هندسی دوربین،
به‌ویژه ماتریس مشخصهٔ دوربین $K$،
بوده است.
در این بخش، رفتار مدل‌ها تحت تغییر این پارامترها
به‌صورت کمی مورد تحلیل قرار می‌گیرد
تا ارتباط عملی میان تحلیل‌های هندسی ارائه‌شده در فصل سوم
و نتایج تجربی برقرار شود.

\subsection{آزمون جاروی $K$: آیا مدل واقعاً فاصلهٔ کانونی را می‌خواند؟}
\label{subsec:intrinsics_variation}

معیارهای تجمیعی مانند AbsRel نشان می‌دهند که ورودی $K$ «سودمند» است، اما به این پرسش سازوکاری پاسخ نمی‌دهند که آیا مدل واقعاً مقدار فاصلهٔ کانونی را می‌خواند و مطابق هندسهٔ تصویربرداری از آن استفاده می‌کند یا خیر. برای پاسخ، آزمونی طراحی شد که در این پژوهش «جاروی $K$»\enfootnote{K-Sweep} نامیده می‌شود: تصویر ورودی ثابت نگه داشته می‌شود و تنها مؤلفه‌های $f_x, f_y$ در ماتریس $K$ ورودی، در ضرایب $\{0.6, 0.8, 1.0, 1.25, 1.6\}$ ضرب می‌شوند. طبق معادلهٔ روزنه‌ای $u = f X / Z$، تصویری یکسان با فاصلهٔ کانونی $\alpha$ برابر، تنها با صحنه‌ای $\alpha$ برابر دورتر سازگار است؛ بنابراین مدلی که $K$ را به‌درستی مصرف می‌کند باید عمق پیش‌بینی‌شدهٔ خود را متناسب با $f$ بازمقیاس کند. سنجهٔ کمی این رفتار، شیب رگرسیون $\log(\text{عمق پیش‌بینی‌شده})$ بر حسب $\log(f)$ است: مقدار ایده‌آل ۱ (پیروی کامل از هندسه) و مقدار ۰ به‌معنای نادیده‌گرفتن کامل $K$ است.

\begin{table}[h]
\centering
\caption{شیب لگاریتمی جاروی $K$ (میانگین $\pm$ انحراف معیار روی ۲۰ تصویر آزمون): پاسخ عمق پیش‌بینی‌شده به مقیاس‌دهی مصنوعی فاصلهٔ کانونی. شیب ایده‌آل برای مدل آگاه از هندسه ۱ است.}
\label{tab:k_sweep_slopes}
\begin{adjustbox}{max width=\textwidth}
\begin{tabular}{lcc}
\hline
زوج آموزشی & \tableproposedmodel{} & \tablecontrolmodel{} \\ \hline
دور ۲ (افزونه‌سازی هندسی) & $0.06 \pm 0.01$ & $0.00$ \\
\textbf{دور ۳ (+ لرزش کانونی)} & $\mathbf{0.93 \pm 0.03}$ & $0.00$ \\ \hline
\end{tabular}
\end{adjustbox}
\end{table}

نتیجهٔ این آزمون (جدول~\ref{tab:k_sweep_slopes} و شکل~\ref{fig:intrinsics_sensitivity}) صریح‌ترین یافتهٔ سازوکاری این پژوهش است. \proposedmodel{} دور ۲ --- با وجود برتری آماری بر \controlmodel{} --- شیب تنها \fadecimal{۰٫۰۶} دارد؛ یعنی عملاً $K$ را نمی‌خواند و برتری‌اش از نشانه‌های غیرمستقیم می‌آید. اما پس از آموزش با لرزش کانونی، شیب به \fadecimal{۰٫۹۳} می‌جهد: خروجی مدل تقریباً دقیقاً همان‌گونه با $f$ مقیاس می‌شود که معادلهٔ $u = f X / Z$ ایجاب می‌کند. جالب آنکه پارامتر دروازهٔ اسکالر کدگذار دوربین در همین مدل همچنان کوچک است ($-0.024$)؛ یعنی اندازهٔ دروازه به‌تنهایی سنجهٔ مناسبی برای میزان مصرف $K$ نیست و بزرگی سیگنال آموخته‌شده در پشت دروازه است که تعیین‌کننده است.

لازم است تأکید شود که تفسیر درست این رفتار برای مدل \emph{متریک}، «پاسخ‌گویی صحیح» به $K$ است، نه «حساسیت نامطلوب»: هنگامی که فاصلهٔ کانونی واقعاً تغییر کند (مثلاً دوربینی با لنز متفاوت)، تصویرِ یکسان به‌واقع با صحنه‌ای در مقیاس متفاوت سازگار است و مدل \emph{باید} پیش‌بینی خود را تغییر دهد؛ \controlmodel{} که به هر $K$ پاسخ یکسان می‌دهد، در چنین شرایطی به‌طور سیستماتیک خطا خواهد کرد. این همان پدیده‌ای است که در ارزیابی میدان دید متغیر زیربخش~\ref{subsec:focal_jitter_results} به‌صورت افت دقت \controlmodel{} مشاهده شد.

\begin{figure}[htbp]
    \centering
    \begin{tikzpicture}[x=1.25cm,y=1.05cm]
    \draw[gray!30, step=1] (0,0) grid (6.6,4.6);
    \draw[->, thick] (0,0) -- (7,0) node[below] {ضریب مقیاس فاصلهٔ کانونی};
    \draw[->, thick] (0,0) -- (0,5) node[left, rotate=90] {عمق نسبی پیش‌بینی‌شده};

    \draw (0,0.08) -- (0,-0.08) node[below] {$0.6$};
    \draw (1.9,0.08) -- (1.9,-0.08) node[below] {$0.8$};
    \draw (3.4,0.08) -- (3.4,-0.08) node[below] {$1.0$};
    \draw (4.9,0.08) -- (4.9,-0.08) node[below] {$1.25$};
    \draw (6.6,0.08) -- (6.6,-0.08) node[below] {$1.6$};
    \draw (0.08,0) -- (-0.08,0) node[left] {$0.6$};
    \draw (0.08,1.2) -- (-0.08,1.2) node[left] {$0.8$};
    \draw (0.08,2.1) -- (-0.08,2.1) node[left] {$1.0$};
    \draw (0.08,3.4) -- (-0.08,3.4) node[left] {$1.25$};
    \draw (0.08,4.6) -- (-0.08,4.6) node[left] {$1.6$};

    \draw[gray, dashed, thick] (0,0) -- (6.6,4.6);
    \draw[blue, thick] plot[mark=*] coordinates {(0,0.25) (1.9,1.18) (3.4,2.1) (4.9,3.47) (6.6,4.6)};
    \draw[teal, thick] plot[mark=triangle*] coordinates {(0,2.05) (1.9,2.07) (3.4,2.1) (4.9,2.18) (6.6,2.35)};
    \draw[red, thick] plot[mark=square] coordinates {(0,2.1) (1.9,2.1) (3.4,2.1) (4.9,2.1) (6.6,2.1)};

    \draw[gray, dashed, thick] (0.3,4.35) -- (0.9,4.35) node[right, black] {ایده‌آل (شیب ۱)};
    \draw[blue, thick] (0.3,4.0) -- (0.9,4.0) node[right, black] {پیشنهادی --- دور ۳};
    \draw[teal, thick] (0.3,3.65) -- (0.9,3.65) node[right, black] {پیشنهادی --- دور ۲};
    \draw[red, thick] (0.3,3.3) -- (0.9,3.3) node[right, black] {کنترل};
    \end{tikzpicture}
    \caption{آزمون جاروی $K$: میانهٔ عمق پیش‌بینی‌شده (نرمال‌شده به حالت ضریب ۱، میانگین هندسی روی ۲۰ تصویر آزمون) در برابر ضریب مقیاس مصنوعی فاصلهٔ کانونی؛ محورها لگاریتمی‌اند. \proposedmodel{} دور ۳ تقریباً به‌طور کامل بر خط ایده‌آل شیب ۱ منطبق است (در ضریب \fadecimal{۱٫۶} دقیقاً عمق \fadecimal{۱٫۶} برابر پیش‌بینی می‌کند)، همان معماری در دور ۲ تقریباً تخت است، و \controlmodel{} طبق تعریف هیچ پاسخی به $K$ ندارد. داده‌ها حاصل اجرای واقعی آزمون بر نقاط وارسی آموزش‌دیده‌اند.}
    \label{fig:intrinsics_sensitivity}
\end{figure}

\subsection{آزمون حذف و تحریف $K$: اثبات اتکای علّی}
\label{subsec:k_knockout}

آزمون جاروی $K$ نشان داد که خروجی مدل \emph{با} $K$ حرکت می‌کند؛ آزمون مکملی لازم است تا نشان دهد دقت مدل نیز \emph{به} $K$ وابسته است و اطلاعات درست دوربین، جزء ضروری عملکرد آن است. برای این منظور، \proposedmodel{} دور ۳ بر کل مجموعهٔ آزمون در چهار شرایط ارزیابی شد: با $K$ صحیح؛ با حذف کامل ورودی $K$؛ و با $K$ عمداً تحریف‌شده ($f_x, f_y$ در \fadecimal{۰٫۷} یا \fadecimal{۱٫۴} ضرب‌شده، تصویر بدون تغییر).

\begin{table}[h]
\centering
\caption{آزمون حذف و تحریف $K$ بر \proposedmodel{} دور ۳ (\lr{ViT-L} متریک، $n=654$، بدون تراز). پیش‌بینی نظری: با شیب جاروی $s=0.93$، تحریف $f$ در ضریب $X$ باید AbsRel را به حدود $|X^{s}-1|$ برساند.}
\label{tab:k_knockout}
\begin{tabular}{lccc}
\hline
شرایط ورودی $K$ & AbsRel$\downarrow$ & $\delta_1\uparrow$ & RMSE$\downarrow$ \\ \hline
$K$ صحیح & $\mathbf{0.0924}$ & $\mathbf{0.9339}$ & $\mathbf{0.3455}$ \\
$K$ حذف‌شده & $0.1467$ & $0.8449$ & $0.4826$ \\
$f \times 0.7$ & $0.2887$ & $0.1310$ & $0.9153$ \\
$f \times 1.4$ & $0.4274$ & $0.1209$ & $1.2118$ \\ \hline
\end{tabular}
\end{table}

نتایج (جدول~\ref{tab:k_knockout}) اتکای علّی را از دو جهت اثبات می‌کند. نخست، حذف $K$ دقت را به‌شدت تخریب می‌کند (AbsRel از \fadecimal{۰٫۰۹۲} به \fadecimal{۰٫۱۴۷}، افت ۵۹ درصدی): مسیر پارامترهای دوربین در این مدل تزئینی نیست و بخش اساسی محاسبهٔ مقیاس از آن عبور می‌کند. دوم --- و از منظر سازوکاری مهم‌تر --- بزرگیِ خطا تحت تحریف $f$ به‌صورت \emph{کمی} با پیش‌بینی هندسی سازگار است: با شیب جاروی \fadecimal{۰٫۹۳}، انتظار می‌رود تحریف $f\times0.7$ خطای مقیاسی حدود $|0.7^{0.93}-1| \approx 0.28$ و تحریف $f\times1.4$ حدود $|1.4^{0.93}-1| \approx 0.37$ ایجاد کند؛ مقادیر مشاهده‌شده (\fadecimal{۰٫۲۸۹} و \fadecimal{۰٫۴۲۷}، شامل خطای پایهٔ مدل) دقیقاً در همین مرتبه‌اند و سقوط $\delta_1$ به حدود \fadecimal{۰٫۱۳} نشان می‌دهد تقریباً تمام پیکسل‌ها به‌اندازهٔ همین ضریب مقیاس جابه‌جا شده‌اند. به بیان دیگر، مدل نه‌تنها از $K$ استفاده می‌کند، بلکه \emph{به همان شکلی} استفاده می‌کند که معادلهٔ تصویربرداری روزنه‌ای تجویز می‌کند --- تصادف یا هم‌بستگی جانبی نمی‌تواند چنین انطباق کمی‌ای تولید کند.

این یافته یک پیامد کاربردی نیز دارد که در بخش محدودیت‌ها (فصل پنجم) بازتاب یافته است: مدلی که به $K$ اتکای علّی دارد، به صحت $K$ نیز حساس است؛ در کاربردهایی که کالیبراسیون نادقیق است، این حساسیت باید در طراحی سامانه لحاظ شود.

\subsection{رفتار متفاوت معیارهای AbsRel و RMSE}
\label{subsec:metric_behavior}

برای تحلیل دقیق‌تر حساسیت مدل‌ها،
از دو معیار AbsRel و RMSE
به‌صورت هم‌زمان استفاده شده است.
نتایج نشان می‌دهد که این دو معیار
رفتارهای متفاوتی نسبت به تغییر پارامترهای دوربین دارند.

معیار AbsRel،
به‌دلیل نرمال‌سازی خطا نسبت به مقدار عمق مرجع،
در بسیاری از موارد تغییرات ناشی از مقیاس
را پنهان می‌کند.
به‌عبارت دیگر،
حتی در شرایطی که خروجی عمق
دچار تغییر مقیاس قابل توجه شده است،
مقدار AbsRel ممکن است
تغییر اندکی را نشان دهد
یا حتی تقریباً ثابت باقی بماند.

در مقابل،
معیار RMSE
به‌طور مستقیم به مقدار مطلق عمق حساس است
و تغییرات مقیاس ناشی از تغییر ماتریس $K$
را به‌وضوح منعکس می‌کند.
نتایج تجربی نشان می‌دهد که
در شرایط تغییر پارامترهای دوربین،
افزایش RMSE در \basemodel{}
به‌مراتب شدیدتر از \proposedmodel{} است،
که این موضوع بیانگر
پایداری بیشتر \proposedmodel{}
از منظر هندسی است.

\subsection{ارتباط با تحلیل هندسی فصل سوم}
\label{subsec:geometry_link}

یافته‌های این بخش
به‌صورت مستقیم با تحلیل‌های هندسی
ارائه‌شده در فصل سوم هم‌راستا هستند.
در فصل سوم نشان داده شد که
عدم لحاظ صریح ماتریس $K$
در فرآیند تخمین عمق
منجر به ابهام مقیاس
و وابستگی خروجی به پارامترهای تصویربرداری می‌شود.

نتایج تجربی این فصل تأیید می‌کنند که
ادغام صریح ماتریس مشخصهٔ دوربین
می‌تواند این ابهام را کاهش دهد
و تخمین عمق را از یک پیش‌بینی نسبی
به تخمینی با مقیاس پایدار نزدیک کند.
بدین ترتیب،
تحلیل حساسیت انجام‌شده در این بخش
نه‌تنها کارایی عملی \proposedmodel{} را نشان می‌دهد،
بلکه اعتبار تجربی
چارچوب هندسی ارائه‌شده در فصل سوم
را نیز تأیید می‌کند.

\section{ملاحظات کیفی و حدود شواهد بصری}
\label{sec:qualitative_results}

بازبینی نقشه‌های عمق در جریان توسعه برای یافتن خروجی صفر، پیکسل‌های مرده و شکست مقیاس به‌کار رفت؛ به‌ویژه همین بازرسی، فروپاشی دو اجرای اولیهٔ \lr{ViT-L} را آشکار کرد. بااین‌حال، شبکهٔ تصویری همسان و قابل بازتولیدی از خروجی تمام نقاط وارسی در مصنوعات نهایی آزمایش‌ها نگهداری نشده است. ازاین‌رو، درج یک شکل نمونه یا ادعای برتری بصری می‌توانست گزینشی و غیرقابل ممیزی باشد و آگاهانه از نسخهٔ نهایی حذف شد.

شاهد اصلی برای تغییر رفتار خروجی، به‌جای قضاوت بصری، آزمون‌های عددی سطح تصویر است: جاروی $K$ پاسخ تقریباً خطی مقیاس پیش‌بینی را نشان می‌دهد، آزمون حذف $K$ افت ۵۹ درصدی ایجاد می‌کند و مقایسهٔ زوجی روی ۶۵۴ تصویر، بهبود AbsRel و RMSE را اندازه‌گیری می‌کند. این شواهد مستقیماً به فرضیهٔ مصرف اطلاعات دوربین مربوط‌اند و از انتخاب چند تصویر مطلوب مصون‌تر هستند.

برای بازتولید کیفی در کار آینده باید شناسهٔ تصاویر پیش از مشاهده قفل شود و برای هر نمونه، تصویر ورودی، عمق مرجع، خروجی \controlmodel{} و خروجی \proposedmodel{} با یک نگاشت رنگ و بازهٔ عمق یکسان ذخیره گردد. تا زمان اجرای چنین پروتکلی، نتیجه‌گیری این پایان‌نامه به شواهد کمی و آزمون‌های علّی گزارش‌شده محدود می‌ماند.


\section{تحلیل موردی دیتاست CityScapes}
\label{sec:cityscapes_analysis}

در میان دیتاست‌های مورد استفاده در این پژوهش،
دیتاست CityScapes
دارای ویژگی‌هایی است که آن را
از سایر مجموعه‌داده‌ها متمایز می‌کند.
به همین دلیل،
رفتار خطاها و نتایج حاصل از ارزیابی مدل‌ها
بر روی این دیتاست
نیازمند تحلیل مستقل و موردی است.

\subsection{ماهیت و ویژگی‌های خاص دیتاست CityScapes}
\label{subsec:cityscapes_nature}

CityScapes
یک دیتاست واقعی شهری\enfootnote{Real-world Urban Dataset}
است که تصاویر آن
در شرایط متنوع نوری، آب‌وهوایی
و با سناریوهای پیچیدهٔ ترافیکی ثبت شده‌اند.
صحنه‌ها شامل اجسام متعدد
در فواصل بسیار متنوع از دوربین هستند،
از خودروها و عابران نزدیک
تا ساختمان‌ها و پس‌زمینه‌های بسیار دوردست.

برخلاف برخی دیتاست‌های دیگر،
CityScapes
به‌صورت ذاتی برای تخمین عمق تک‌دوربینه طراحی نشده
و عمق مرجع آن
از طریق تطبیق استریو و بازسازی سه‌بعدی
به‌دست آمده است.
این موضوع منجر به تفاوت‌های ساختاری
در کیفیت و یکنواختی داده‌های عمق مرجع می‌شود.

\subsection{تحلیل کیفیت عمق مرجع}
\label{subsec:cityscapes_gt}

عمق مرجع در CityScapes
دارای چگالی و دقت یکنواخت در سراسر تصویر نیست.
در نواحی نزدیک به دوربین
و بر روی اجسام برجسته،
عمق مرجع از کیفیت نسبتاً بالایی برخوردار است،
اما در نواحی دوردست،
لبه‌ها و سطوح با بافت ضعیف،
خطاها و ناپیوستگی‌های قابل توجهی مشاهده می‌شود.

این ناهمگنی در کیفیت عمق مرجع
باعث می‌شود که
معیارهای متریک،
به‌ویژه RMSE و SiLog،
نسبت به نویز و خطاهای محلی
حساسیت بیشتری نشان دهند.
در نتیجه،
مقادیر خطای گزارش‌شده
لزومأ بازتاب‌دهندهٔ صرف
عملکرد واقعی مدل نیستند،
بلکه تا حدی تحت تأثیر محدودیت‌های دادهٔ مرجع قرار دارند.

\subsection{دلیل رفتار متفاوت خطاها}
\label{subsec:cityscapes_error_behavior}

نتایج تجربی نشان می‌دهد که
رفتار خطاها در CityScapes
با برخی دیتاست‌های دیگر
مانند \lr{NYUv2} یا VKITTI
تفاوت معناداری دارد.
یکی از دلایل اصلی این تفاوت،
دامنهٔ بسیار وسیع عمق در صحنه‌های شهری است،
که از چند متر تا ده‌ها متر متغیر است.

در چنین شرایطی،
خطاهای کوچک نسبی
می‌توانند به خطاهای مطلق بزرگ
در معیارهای متریک منجر شوند.
از سوی دیگر،
وجود نواحی دوردست با عمق نامطمئن
باعث می‌شود که
مدل‌هایی با رفتار پایدارتر از نظر مقیاس
الزاماً بهبود یکنواختی
در تمامی معیارها نشان ندهند.

\subsection{محدودیت همسان‌سازی نتایج با سایر دیتاست‌ها}
\label{subsec:cityscapes_limitations}

با توجه به ویژگی‌های ذکرشده،
همسان‌سازی مستقیم مقادیر خطا
بین CityScapes و سایر دیتاست‌ها
دارای محدودیت‌های ذاتی است.
تفاوت در نحوهٔ تولید عمق مرجع،
دامنهٔ عمق صحنه،
و کیفیت داده‌ها
باعث می‌شود که
مقایسهٔ عددی ساده
بدون در نظر گرفتن این عوامل
می‌تواند منجر به برداشت نادرست شود.

از این رو،
در این پژوهش،
نتایج CityScapes
بیشتر به‌عنوان یک مطالعهٔ موردی
برای بررسی رفتار مدل‌ها
در شرایط شهری واقعی و پیچیده
مورد استفاده قرار گرفته است،
نه به‌عنوان مرجع مطلق
برای مقایسهٔ کمی با سایر دیتاست‌ها.
این رویکرد امکان تحلیل واقع‌بینانه‌تر
و تفسیر دقیق‌تر نتایج
را فراهم می‌کند.


\section{بحث و تفسیر نتایج}
\label{sec:discussion}

در این بخش، نتایج کمی و کیفی ارائه‌شده در فصل حاضر
به‌صورت یکپارچه مورد بحث و تفسیر قرار می‌گیرند.
هدف اصلی این بخش،
فراتر رفتن از گزارش صرف اعداد و نمودارها
و پاسخ‌گویی صریح به سوالات تحقیق،
تبیین رفتار مدل‌ها،
و تفسیر یافته‌ها در چارچوب هندسهٔ تصویربرداری
و ادبیات موجود در حوزهٔ تخمین عمق تک‌دوربینه است.

\subsection{جمع‌بندی تحلیلی نتایج کمی و کیفی}
\label{subsec:quantitative_qualitative_summary}

نتایج کمی نشان دادند که
\basemodel{}،
اگرچه در معیارهای غیرمتریک
عملکرد رقابتی و قابل قبولی دارد،
اما در معیارهای متریک،
به‌ویژه SiLog و RMSE،
رفتاری ناپایدار نسبت به تغییر دیتاست
و پارامترهای دوربین از خود نشان می‌دهد.
این ناپایداری به‌صورت اختلاف‌های معنادار آماری
در مقایسهٔ بین‌دیتاستی آشکار شد.

در مقابل،
\proposedmodel{} که در آن
ماتریس مشخصهٔ دوربین به‌صورت صریح لحاظ شده است،
بهبود قابل توجهی در معیارهای متریک
و کاهش معنادار حساسیت نسبت به تغییر دیتاست‌ها
نشان داد.
این بهبودها نه‌تنها از نظر آماری معنادار بودند،
بلکه در تحلیل‌های کیفی نیز
به‌صورت افزایش یکنواختی سطوح،
پایداری مقیاس اجسام
و بازنمایی واقع‌بینانه‌تر نواحی دوردست مشاهده شدند.

هم‌راستایی نتایج کمی و کیفی
نشان می‌دهد که
بهبودهای حاصل‌شده
صرفاً نتیجهٔ بهینه‌سازی عددی معیارها نیستند،
بلکه بیانگر تغییر معنادار
در نحوهٔ درک مدل از ساختار هندسی صحنه هستند.

\subsection{پاسخ به سوالات تحقیق}
\label{subsec:research_questions_answers}

یافته‌های این پژوهش امکان پاسخ‌گویی مستقیم
به سوالات تحقیق مطرح‌شده در فصل اول را فراهم می‌کند.

نخست،
در پاسخ به این سوال که
آیا می‌توان با تکیه بر یک مدل بنیادین تک‌دوربینه
به عمق متریک پایدار دست یافت،
نتایج نشان می‌دهد که فاین‌تیون نظارتی \basemodel{}
دقت متریک را به‌طور چشمگیری بهبود می‌دهد؛ بااین‌حال،
وقتی مقیاس در داده عمداً مبهم شود، مدل فاقد $K$
با یک کف خطای کاهش‌ناپذیر مواجه می‌شود. بنابراین،
اطلاعات تصویر می‌تواند مقیاس تقریبی را از نشانه‌های معنایی بیاموزد،
اما در حالت هندسی مبهم تضمینی برای بازیابی یکتای مقیاس ندارد.

دوم،
در خصوص نقش ماتریس مشخصهٔ دوربین،
نتایج به‌وضوح نشان می‌دهند که
ادغام صریح $K$
می‌تواند وابستگی خروجی عمق به پارامترهای تصویربرداری
را کاهش داده
و پایداری متریک مدل را به‌طور معناداری افزایش دهد.

در نهایت،
پژوهش حاضر نشان می‌دهد که
بهبود معیارهای متریک
نه از طریق تغییرات پیچیده در معماری،
بلکه با افزودن اطلاعات هندسی مناسب
و هم‌راستا با فیزیک تصویربرداری
قابل دستیابی است.

\subsection{وابستگی سهم پارامترهای دوربین به پروتکل ارزیابی و ظرفیت مدل}
\label{subsec:protocol_capacity_dependence}

یافته‌های مطالعهٔ کنترل‌شدهٔ بخش~\ref{sec:controlled_nyu_study}
تصویری دقیق‌تر و مشروط‌تر از سهم پارامترهای درونی دوربین ترسیم می‌کنند
که شایستهٔ بحث مستقل است.

نخست، سهم این پارامترها به \emph{پروتکل ارزیابی} وابسته است.
در پروتکل نسبی، تراز مقیاس و جابجایی به‌ازای هر تصویر،
دقیقاً همان مؤلفهٔ اطلاعاتی را از ارزیابی حذف می‌کند
که میدان دید و فاصلهٔ کانونی در اختیار مدل می‌گذارند؛
یعنی مقیاس مطلق صحنه.
از این رو، مشاهدهٔ برتری‌نیافتن \proposedmodel{} بر \controlmodel{}
در پروتکل نسبی \lr{ViT-L}
نه شاهدی علیه فرضیهٔ پژوهش،
بلکه پیامد ساختاری سازوکار ارزیابی است.
در مقابل، در ارزیابی متریک بدون تراز،
برتری \proposedmodel{} در تمامی معیارها ظاهر شده
و در RMSE در هر دو دور آزمایش معنادار است.
مهم‌تر آنکه الگوی پاسخ به دوز
--- رشد فاصلهٔ دو مدل با افزایش تنوع میدان دید در داده ---
نشان می‌دهد که این برتری، پیوند علّی با اطلاعات هندسی دارد
و صرفاً نوسان آماری نیست.

دوم، سهم پارامترهای دوربین به \emph{ظرفیت انکودر} نیز وابسته است.
در اندازهٔ \lr{ViT-S}، حتی در پروتکل نسبی،
افزودن $K$ بهبود کوچک اما معنادار ایجاد کرد؛
در حالی که در \lr{ViT-L} چنین نبود.
تبیین محتمل آن است که انکودر بزرگ‌تر،
سرنخ‌های پرسپکتیوی و معنایی تصویر
(نظیر اندازهٔ ظاهری اجسام آشنا و هندسهٔ خطوط گریز)
را با دقتی استخراج می‌کند که برای بازسازی ساختار نسبی،
اطلاعات مکمل چندانی از $K$ باقی نمی‌ماند؛
حال آنکه انکودر کوچک‌تر از هر سیگنال کمکی بهره می‌برد.
این مشاهده با یافته‌های مربوط به جبران میدان دید در ادبیات
\cite{Saxena2023FOV}
هم‌راستا است و مرزهای سودمندی رویکرد پیشنهادی را
به‌صورت تجربی مشخص می‌کند.

سوم، نتایج نشان دادند که بخش قابل توجهی از بهبود مطلق
نسبت به \publishedbasemodel{}،
حاصل ترکیب فاین‌تیون درون‌دامنه‌ای بدون نشت
و افزونه‌سازی هندسی صحیح است
و سهم خالص ورودی $K$، افزایشی و در مرتبهٔ کوچک‌تری قرار دارد.
گزارش شفاف این تفکیک،
برای جلوگیری از بیش‌برآورد سهم هر مؤلفه ضروری است
و اعتبار علمی ادعاهای پژوهش را تضمین می‌کند.

\subsection{مقایسهٔ غیرمستقیم با ادبیات موضوع}
\label{subsec:literature_comparison}

در بسیاری از مطالعات پیشین
در حوزهٔ تخمین عمق تک‌دوربینه،
تمرکز اصلی بر بهبود معیارهای غیرمتریک
یا استفاده از نرمال‌سازی‌های ضمنی
برای کاهش اثر مقیاس بوده است.
در این رویکردها،
اطلاعات دوربین یا نادیده گرفته شده
یا به‌صورت ضمنی در داده‌ها نهفته فرض شده است.

در مقابل،
یافته‌های این پژوهش نشان می‌دهد که
نادیده گرفتن صریح هندسهٔ دوربین
می‌تواند منجر به ارزیابی‌های گمراه‌کننده
در معیارهای نرمال‌شده شود،
در حالی که خطاهای متریک واقعی پنهان باقی می‌مانند.
از این منظر،
نتایج حاضر هم‌راستا با دیدگاه‌هایی است
که بر لزوم بازگشت به مدل‌های فیزیکی تصویربرداری
در طراحی شبکه‌های یادگیری عمیق تأکید دارند،
هرچند مسیر پیاده‌سازی آن‌ها
در این پژوهش متفاوت و ساده‌تر دنبال شده است.

\subsection{تفسیر نتایج در چارچوب هندسهٔ تصویربرداری}
\label{subsec:geometric_interpretation}

از منظر هندسهٔ تصویربرداری،
عمق متریک به‌صورت مستقیم
به پارامترهای کانونی و هندسهٔ پروجکشن وابسته است.
بنابراین،
مدلی که تنها بر اطلاعات تصویری تکیه دارد،
در بهترین حالت
می‌تواند یک بازنمایی نسبی از عمق را بیاموزد.

نتایج این پژوهش تأیید می‌کنند که
ادغام صریح ماتریس مشخصهٔ دوربین
در فرآیند تخمین عمق،
در واقع محدودیت‌های فیزیکی سیستم تصویربرداری
را به فضای یادگیری تحمیل می‌کند.
این محدودیت‌ها باعث می‌شوند که
مدل به‌جای اتکا به هم‌بستگی‌های آماری ناپایدار،
وابستگی‌های هندسی معنادار
میان تصویر و عمق را یاد بگیرد.

در نتیجه،
بهبود پایداری متریک مشاهده‌شده
را می‌توان به‌عنوان
پیامد مستقیم هم‌راستا شدن
معماری یادگیری با اصول هندسی
مدل سوراخ‌سوزنی تفسیر کرد.
این تفسیر نشان می‌دهد که
ترکیب آگاهانهٔ یادگیری عمیق
و هندسهٔ کلاسیک
می‌تواند مسیر مؤثری
برای پیشرفت در تخمین عمق تک‌دوربینه فراهم کند.


\section{جمع‌بندی فصل نتایج}
\label{sec:results_summary}

در این فصل، نتایج حاصل از ارزیابی کمی و کیفی
\proposedmodel{} و \basemodel{}
در چارچوبی منسجم ارائه و تحلیل شد.
هدف اصلی این فصل،
پاسخ‌گویی تجربی به سوالات تحقیق
و بررسی میزان تحقق اهداف پژوهش
بر اساس داده‌ها و آزمایش‌های انجام‌شده بود.

نتایج کمی نشان دادند که
\basemodel{}،
علی‌رغم عملکرد مناسب در معیارهای غیرمتریک،
در معیارهای متریک رفتاری ناپایدار
نسبت به تغییر دیتاست و پارامترهای دوربین دارد.
این ناپایداری به‌صورت اختلاف‌های معنادار آماری
در مقایسه‌های بین‌دیتاستی
و حساسیت بالا در معیارهایی مانند RMSE و SiLog
مشاهده شد.

در مقابل،
\proposedmodel{} که در آن
ماتریس مشخصهٔ دوربین به‌صورت صریح
در فرآیند تخمین عمق لحاظ شده است،
بهبود معناداری در معیارهای متریک نشان داد
و حساسیت آن نسبت به تغییر شرایط تصویربرداری
به‌طور محسوسی کاهش یافت.
این بهبودها نه‌تنها از نظر آماری معنادار بودند،
بلکه در تحلیل‌های کیفی نیز
به‌صورت پایداری مقیاس،
یکنواختی بهتر سطوح
و بازنمایی واقع‌بینانه‌تر نواحی دوردست
قابل مشاهده بودند.

تحلیل حساسیت نسبت به پارامترهای دوربین
نشان داد که
معیارهای نرمال‌شده‌ای مانند AbsRel
توانایی محدودی در آشکارسازی خطاهای مقیاس دارند،
در حالی که معیارهای متریک
رفتار هندسی واقعی مدل‌ها را بهتر منعکس می‌کنند.
این یافته به‌طور مستقیم
تحلیل‌های نظری فصل سوم
در خصوص نقش ماتریس $K$
در بازیابی عمق متریک را تأیید می‌کند.

همچنین،
تحلیل موردی دیتاست CityScapes
نشان داد که
تفاوت در ماهیت داده‌ها
و کیفیت عمق مرجع
می‌تواند بر رفتار معیارهای ارزیابی اثرگذار باشد
و تفسیر نتایج نیازمند در نظر گرفتن
محدودیت‌های ذاتی هر دیتاست است.

مطالعهٔ کنترل‌شدهٔ بخش~\ref{sec:controlled_nyu_study}،
که پس از اصلاح تفکیک آموزش/آزمون \lr{NYUv2}
و راستی‌آزمایی زنجیرهٔ ارزیابی از طریق بازتولید نتایج مقالهٔ مرجع انجام شد،
دقیق‌ترین سنجش این پژوهش از سهم پارامترهای درونی دوربین را فراهم آورد.
بر اساس این مطالعه،
\proposedmodel{} \lr{ViT-L} در پروتکل استاندارد،
از \publishedbasemodel{}
در تمامی معیارها به‌صورت معنادار آماری بهتر است
(کاهش AbsRel از \fadecimal{۰٫۰۴۲۵} به \fadecimal{۰٫۰۳۷۷} در پروتکل نسبی
و از \fadecimal{۰٫۲۰۷} به \fadecimal{۰٫۰۷۹} در ارزیابی متریک).
سهم خالص ورودی $K$ در برابر \controlmodel{} هم‌داده،
در پروتکل نسبی تنها برای انکودر کوچک معنادار بود،
اما در ارزیابی متریک بدون تراز،
در هر دو دور آزمایش و به‌صورت فزاینده با شدت تنوع میدان دید
(الگوی پاسخ به دوز) آشکار شد.
این نتایج، همراه با تحلیل توضیحی نقش سازوکار تراز در پنهان‌سازی اطلاعات مقیاس،
درک دقیق‌تری از شرایط سودمندی اطلاعات هندسی دوربین به‌دست می‌دهند.

در مجموع،
نتایج این فصل نشان می‌دهد که
ادغام آگاهانهٔ اطلاعات هندسی دوربین
در مدل‌های تخمین عمق تک‌دوربینه
می‌تواند گامی مؤثر
در جهت دستیابی به عمق متریک پایدار باشد.
این فصل، با ارائهٔ شواهد تجربی،
زمینه را برای جمع‌بندی نهایی پژوهش
و بیان دستاوردها، محدودیت‌ها
و مسیرهای آیندهٔ تحقیق
در فصل بعد فراهم می‌کند.
